\documentclass[11pt]{article}
\usepackage{iftex}
\ifPDFTeX
  \usepackage[T1]{fontenc}
  \usepackage{lmodern}
\fi
\usepackage{amsmath,amssymb,amsthm}
\usepackage[margin=1in]{geometry}
\usepackage{booktabs,longtable,array}
\usepackage[numbers,sort&compress]{natbib}
\usepackage{microtype}
\usepackage{xurl}
\usepackage{hyperref}
\hypersetup{colorlinks=true,linkcolor=blue,citecolor=blue,urlcolor=blue,
  pdftitle={Rough numbers at a fixed distance with prescribed Liouville signs},
  pdfauthor={Jizhou Guo}}
\Urlmuskip=0mu plus 1mu
\setlength{\emergencystretch}{2em}
\newtheorem{theorem}{Theorem}[section]
\newtheorem{lemma}[theorem]{Lemma}
\newtheorem{proposition}[theorem]{Proposition}
\newtheorem{corollary}[theorem]{Corollary}
\newtheorem{inputthm}[theorem]{Input}
\theoremstyle{definition}
\newtheorem{definition}[theorem]{Definition}
\newcommand{\e}{\mathrm e}
\newcommand{\E}{\mathbb E}
\newcommand{\Prob}{\mathbb P}
\newcommand{\Z}{\mathbb Z}
\newcommand{\N}{\mathbb N}
\newcommand{\C}{\mathbb C}
\newcommand{\R}{\mathbb R}
\newcommand{\ind}{\mathbf 1}
\newcommand{\Dpret}{\mathbb D}
\newcommand{\ph}{\varphi}
\newcommand{\leanname}[1]{\nolinkurl{#1}}
\newcommand{\norm}[1]{\left\lVert#1\right\rVert}
\newcommand{\abs}[1]{\left\lvert#1\right\rvert}
\DeclareMathOperator{\rad}{rad}
\DeclareMathOperator{\tr}{tr}
\DeclareMathOperator{\HS}{HS}
\DeclareMathOperator{\Rea}{Re}
\DeclareMathOperator{\cond}{cond}
\title{Rough numbers at a fixed distance with prescribed Liouville signs}
\author{Jizhou Guo\\
  \href{mailto:mitsuha2021b@gmail.com}{\texttt{mitsuha2021b@gmail.com}}}
\date{October 2026}

\DeclareMathOperator{\supp}{supp}
\newcommand{\Pset}{\mathcal P}
\newcommand{\Sset}{\mathcal S}
\newcommand{\Rz}{R_z}
\newcommand{\wt}{W_{\mathbf t}}
\theoremstyle{remark}
\newtheorem{remark}[theorem]{Remark}

\numberwithin{equation}{section}
\newcommand{\Qset}{\mathcal Q}
\newcommand{\Dset}{\mathcal D}
\newcommand{\eps}{\varepsilon}
\newcommand{\sfw}{\mathrm{sf}}
\newcommand{\cprod}[1]{\prod_{p\mid #1}(\ind_{p\mid n}-1/p)}
\begin{document}
\maketitle
\begin{abstract}
Let $\lambda$ be the Liouville function and let $h\ge2$ be a fixed even integer. We show that there is an absolute
constant $\gamma>0$ such that, uniformly for $2\le z\le\exp((\log X)^{1/3750})$, each of the four sign patterns of
$(\lambda(n),\lambda(n+h))$ occurs for $\tfrac14XV_h(z)\,(1+O_h((\log X)^{-\gamma}))$ integers $n\le X$ such that
$n$ and $n+h$ have no prime factor up to $z$; here $XV_h(z)$ is the main term for the number of such $n$.
Consequently each pattern occurs for $\gg_hX/(\log X)^{1/1875}$ integers $n\le X$ with $n$ and $n+h$ squarefree,
free of prime factors up to $\exp((\log X)^{1/3750})$, and with at most $(1-1/4000)\log\log X$ prime factors. We
prove corresponding statements for the weights $r^{\omega_z(n)}$ on the small prime factors and
$r^{\Omega(n)}$ on all prime factors, in each case with a squarefreeness restriction at the small primes. The proofs adapt the graph argument of OpenAI's proof of the two-point Chowla
conjecture, keeping the sieve density inside the operator.
\end{abstract}
\noindent\textit{2020 Mathematics Subject Classification.} Primary 11N37; Secondary 11N36, 11N25, 11P32.

\section{Introduction}\label{sec:intro}
All logarithms are natural. The function $\Omega(n)$ counts prime factors with multiplicity, $\omega(n)$ counts
distinct prime factors, and $\lambda(n)=(-1)^{\Omega(n)}$. Put $P^-(1)=\infty$ and write $P^-(n)$ for the least
prime factor otherwise. Throughout the paper $h\ge2$ is a fixed even integer. Define
\begin{equation}\label{eq:rough-def}
 R_z(n)=\ind_{P^-(n)>z},\qquad
 V_h(z)=\prod_{p\le z}\left(1-\frac{\nu_h(p)}p\right),\qquad
 \nu_h(p)=\begin{cases}1,&p\mid h,\\2,&p\nmid h.\end{cases}
\end{equation}
Sums and counts over $n\le X$ always concern positive integers.

\begin{theorem}\label{thm:rough}
There is an absolute constant $\gamma>0$ such that, for each fixed even $h\ge2$ and every sufficiently large $X$,
uniformly for $2\le z\le\exp((\log X)^{1/3750})$ and $(s,t)\in\{\pm1\}^2$,
\begin{equation}\label{eq:rough-main}
\#\{n\le X:R_z(n)R_z(n+h)=1,\ (\lambda(n),\lambda(n+h))=(s,t)\}
 =\tfrac14XV_h(z)\bigl(1+O_h((\log X)^{-\gamma})\bigr).
\end{equation}
\end{theorem}

\begin{theorem}\label{thm:omega-consumer}
For each fixed even $h\ge2$ and each sign pattern there are $\gg_h X/(\log X)^{1/1875}$ integers $n\le X$ such that
$n$ and $n+h$ are squarefree, have no prime factor up to $\exp((\log X)^{1/3750})$, have that sign pattern, and satisfy
\begin{equation}\label{eq:hard-omega}
 \Omega(n),\Omega(n+h)\le(1-1/4000)\log\log X.
\end{equation}
This holds for every sufficiently large $X$.
\end{theorem}

For $z\ge2$ let
\begin{equation}\label{eq:soft-def}
 S_z(n)=\prod_{p\le z}\ind_{p^2\nmid n},\qquad
 \omega_z(n)=\#\{p\le z:p\mid n\},\qquad b_z(n)=S_z(n)r^{\omega_z(n)}.
\end{equation}
\begin{theorem}\label{thm:soft}
There is an absolute constant $\gamma>0$, which may be the same as in Theorem~\ref{thm:rough}, such that for each
fixed even $h\ge2$, each fixed $0<r<1$, and every sufficiently large $X$, uniformly for
$2\le z\le\exp((\log X)^{1/3750})$ and $(s,t)\in\{\pm1\}^2$,
\begin{equation}\label{eq:soft-main}
 \sum_{\substack{n\le X\\(\lambda(n),\lambda(n+h))=(s,t)}}b_z(n)b_z(n+h)
 =\tfrac14X V^{\sfw}_{r,h}(z)\bigl(1+O_{h,r}((\log X)^{-\gamma})\bigr),
\end{equation}
where $V^{\sfw}_{r,h}(z)=\prod_{p\le z}v_{r,h}(p)$ and
\begin{equation}\label{eq:soft-local}
 v_{r,h}(p)=\begin{cases}
 1-2/p+2r(p-1)/p^2,&p\nmid h,\\
 1-1/p+r^2(p-1)/p^2,&p^2\mid h,\\
 1-1/p+r^2(p-2)/p^2,&p\mid h,\ p^2\nmid h.
 \end{cases}
\end{equation}
The threshold and implied constant may depend on $h,r$; $r$ is fixed before $X$.
\end{theorem}

The constants for the following theorem are separate from those used in Theorems~\ref{thm:rough} and~\ref{thm:soft}.
Let $C_3$ be the absolute graph constant in \cite[(8.7), p.~41]{OpenAI007}, with the same permissible choice for the
two-layer operators of Proposition~\ref{prop:phase}. Put $C_\Omega=\max(4,C_3)$ and fix
\begin{equation}\label{eq:full-fixed}
 W_\Omega=16\e C_\Omega^2,\quad \alpha_\Omega=\frac1{1200W_\Omega},\quad
 \delta_\Omega=\sqrt{\alpha_\Omega/216},\quad r=1-\delta_\Omega,\quad
 \gamma_\Omega=\alpha_\Omega/3600.
\end{equation}
\begin{theorem}\label{thm:full}
With the parameters \eqref{eq:full-fixed} and $B(n)=r^{\Omega(n)}S_y(n)$, for each fixed even $h\ge2$, every
sufficiently large $X$, $2\le y\le\exp((\log X)^{1/216})$, and every nonzero
$(e_1,e_2)\in\{0,1\}^2$,
\begin{equation}\label{eq:full-main}
 \left|\sum_{n\le X}\lambda(n)^{e_1}\lambda(n+h)^{e_2}B(n)B(n+h)\right|
 \ll_h X(\log X)^{-2\delta_\Omega-\gamma_\Omega}.
\end{equation}
In the same range the unsigned mass satisfies
\begin{equation}\label{eq:full-mass-main}
 \sum_{n\le X}B(n)B(n+h)\asymp_h X(\log X)^{-2\delta_\Omega}.
\end{equation}
The parameter $y$ imposes squarefreeness at small primes; it imposes no roughness restriction.
\end{theorem}

\begin{remark}\label{rem:tradeoff}
Put $\rho=1/3750$. For every fixed $0<\kappa\le\rho$ and $0<\epsilon<\kappa$, for every sufficiently large $X$, every sign pattern occurs for at least
$c_{h,\kappa,\epsilon}X/(\log X)^{2\kappa}$ integers $n\le X$ with both endpoints squarefree and
\begin{equation}\label{eq:tradeoff}
 \Omega(n),\Omega(n+h)\le(1-\kappa+\epsilon)\log\log X.
\end{equation}
For $\kappa=\rho$ both endpoints may additionally be required to have no prime factor up to
$\exp((\log X)^\rho)$. For $\kappa<\rho$ this additional condition is not part of the assertion.
All parameters are fixed before $X$.
\end{remark}

Mertens' theorem gives $V_h(z)\asymp_h(\log z)^{-2}$; see Lemma~\ref{lem:density}.
At the upper endpoint of Theorem~\ref{thm:rough}, the density is of order $(\log X)^{-2/3750}$.
An absolute cancellation estimate must therefore save more than this density to give a relative asymptotic.
The bound on $\Omega$ in Theorem~\ref{thm:omega-consumer} grows with $X$.

OpenAI's paper of 24 September 2026 \cite[Theorem~1.1 and Proposition~2.1, pp.~3, 7]{OpenAI007} gives an absolute
power-of-logarithm saving for Liouville correlations in fixed affine forms and fixed progressions, including
nonunit classes. For a fixed $z$, finite inclusion--exclusion over the primes up to $z$, together with classical
one-point cancellation, gives each rough sign pattern a quarter of the rough-pair mass. Its
Theorem~1.2 (p.~4) gives qualitative ordinary cancellation for bounded multiplicative functions when an original
factor is uniformly nonpretentious, also with fixed progression restrictions.

Goldston, Graham, Pintz and Y\i ld\i r\i m (18 March 2008) \cite[Theorems~6 and~9, pp.~5, 7]{GGPPY} prove that
$\Omega(n)=\Omega(n+2)=B$ infinitely often for each integer $B\ge5$, and for each fixed even $h$ that
$\Omega(n)=\Omega(n+h)=A$ infinitely often when $A\ge4$ if $4\mid h$, or $A\ge5$ if $h\equiv2\pmod4$.
Helfgott and Radziwi\l\l\ (11 March 2021; version~2, 13 April 2021) \cite[Corollaries~1.4, 1.7--1.8, pp.~5--7]{HR}
give correlations with restrictions on $\Omega$ in logarithmic averages and at almost all scales.
Matom\"aki, Radziwi\l\l\ and Tao (4 September 2015; version~2, 23 September 2015)
\cite[Theorems~1.6 and~1.9, pp.~2--3]{MRTpatterns} prove positive lower natural density for all eight Liouville
patterns of length three and all nine adjacent M\"obius patterns.

Weighted two-point sums of the present kind were formulated as a conjecture by Preobrazhenski\u{\i} and
Preobrazhenskaya \cite[Conjecture~2, p.~3]{PP2015}. For $y\ge2$ let $\omega^-(n,y)$ and $\omega^+(n,y)$ be the
numbers of prime divisors of $n$ that are less than $y$ and at least $y$, and on squarefree $n$ put
$\tau^{\pm}(n)=\tau^{\pm}_{\kappa_1,\kappa_2}(n,y)=\kappa_1^{\omega^-(n,y)}\kappa_2^{\omega^+(n,y)}$.
Their conjecture states that, for fixed $0\le\kappa_1<\kappa_2$ and $y=\exp((\log x)^\delta)$ with some fixed
$0<\delta<1$,
\[
\sum_{N/2\le n\le N}\mu(n+h_1)\mu(n+h_2)\tau^{\pm}(n+h_1)\tau^{\pm}(n+h_2)
=o\Bigl(\sum_{N/2\le n\le N}\tau^{\pm}(n+h_1)\tau^{\pm}(n+h_2)\Bigr).
\]
They remark that the case $\kappa_1=\kappa_2=1$ is the two-point Chowla conjecture for $\mu$, and state that for
$0\le\kappa_1<\kappa_2$ the conjecture essentially implies the twin prime conjecture; their Theorem~4 includes a
logarithmically averaged two-point estimate for $\lambda$ with a sieve weight. The weights $R_z$ and $b_z$ of
Theorems~\ref{thm:rough} and~\ref{thm:soft} have this shape with $\kappa_2=1$ and $\kappa_1=0$, respectively
$\kappa_1=r$. The two theorems are stated for $\lambda$, impose squarefreeness only at the primes up to $z$, and
hold for $z\le\exp((\log X)^{1/3750})$.

For the truncated functions $\lambda_y(n)=(-1)^{\Omega_y(n)}$ and $\mu_y(n)=\mu^2(n)(-1)^{\omega_y(n)}$, where
$\Omega_y$ and $\omega_y$ count only the prime factors up to $y$, Daboussi and S\'ark\"ozy
\cite[Corollary~2]{DaboussiSarkozy2003} prove
$|\sum_{n\le x}\lambda_y(n)\lambda_y(n+1)|\ll x((\log y)^{-9}+\exp(-u/8))$ for $3\le y\le x$ and $u=\log x/\log y$,
and Mangerel \cite[Theorem~1.2(a)]{Mangerel2017} proves, for each fixed $a$ and in a range of
$\beta=\log x/\log y$ specified there,
$\sum_{n\le x}\mu_y(n)\mu_y(n+a)\ll x(1/\log\log y+\exp(-\tfrac1{21}\beta\log\beta))$. Mangerel's introduction
attributes the first estimate of this type, with $y\le(\log x)^{2+o(1)}$, to Cassaigne, Ferenczi, Mauduit, Rivat
and S\'ark\"ozy \cite{CFMRS1999}. These results concern the prime factors up to $y$;
Theorem~\ref{thm:rough} concerns integers without such prime factors.

Without sieve conditions, Tao \cite[Theorem~1.2]{Tao2016} proved the logarithmically averaged two-point Chowla
conjecture, Tao and Ter\"av\"ainen \cite[Corollary~1.14]{TaoTeravainen2019} proved that
$\sum_{n\le X}\lambda(n+h_1)\lambda(n+h_2)=o(X)$ for $X$ outside a set of logarithmic density zero, and Pilatte
\cite[Theorem~1.1]{Pilatte2023} proved $\sum_{n\le x}\lambda(n)\lambda(n+1)/n\ll(\log x)^{1-c}$ with an absolute
$c>0$. Tao and Ter\"av\"ainen \cite[Theorem~3.1]{TaoTeravainen2025} give a quantitative two-point estimate for
$1$-bounded multiplicative functions outside an exceptional set of scales, and deduce an asymptotic formula for
consecutive smooth numbers at those scales \cite[Theorem~1.8]{TaoTeravainen2025}; a second paper of the release
containing [O] proves the joint Dickman law for the largest prime factors of $n$ and $n+1$ in natural density
\cite[Theorem~1.1]{OpenAI012}. Charamaras and Richter \cite[Theorem~A]{CharamarasRichter2024} prove asymptotic
independence of $\Omega(n)$ and $\Omega(n+1)$ along logarithmic averages, and Hughes
\cite[Corollary~7.1]{Hughes2026} gives logarithmic frequency $1/4$ for each sign pair of $\lambda$ at two affine
forms in every window $(x/\omega,x]$ with $\omega$ above an explicit threshold.

A Lean~4 repository \cite{JyhSalt2026} contains formal statements at specific small sieve levels: for every $z$
with $\prod_{p\le z}p\le548$ (that is, $z\le10$) and every admissible residue class modulo that product,
infinitely many $n$ with $n(n+2)$ free of prime factors up to $z$ and $\Omega(n(n+2))$ odd (5 September 2026); and,
when the product is at most $2310$, a lower bound for the logarithmic mass of such $n$ in each admissible class
(14 September 2026), described there as an instance of \cite{Tao2016}. The sifted sum
$\sum_{n\le N,\ P^-(n)>D,\ P^-(n+2)>D}\lambda(n)\lambda(n+2)$ with $D=N^{1/8}$ is considered in
\cite[Section~9]{McCaffer2026}, in a reduction of the twin prime conjecture to a hypothesis on
$\lambda(n)\lambda(n+2)$ in arithmetic progressions.

The author's papers \cite{GuoRough,GuoShifted}, both dated 9 October 2026, treat, respectively, sums of two rough
numbers with prescribed Liouville signs and shifted primes with a prescribed parity of the number of prime
factors. The paper \cite{GuoEarlier}, dated 8 October 2026, treats Goldbach-type representations with prescribed
Liouville signs. The present argument concerns a fixed distance between the two integers.

\paragraph{What is used.}
Write [O] for \cite{OpenAI007}, taken at repository commit
\begin{center}\small\nolinkurl{adc7f1241b42e322a6451854ab7e4b4c146bf78a}.\end{center}
We use [O, Lemma~5.5, p.~26], the arithmetic and combinatorial statements
[O, Lemmas~6.1, 6.3--6.4, 7.1--7.2, pp.~27, 30--37], and the finite-dimensional transfer
[O, Lemma~8.1, p.~39] as stated, checking their hypotheses in Proposition~\ref{prop:masked-moment}.
For the padding weight $16^{\omega(u)}$, vertex weight $17^{\omega_{\Qset}}$, rough-target cutoff, balanced soft
weight and numerical signs in the full-$\Omega$ argument, the original statements of [O, Lemma~2.2,
Lemmas~4.1--4.2, 4.4--4.6, 5.2--5.3, 5.6, 6.2, 6.5, 7.3, Theorem~5.4, Lemma~8.2 and Propositions~8.3--8.4] do not
serve as direct instances. We give the required versions in Sections~\ref{sec:setup}--\ref{sec:rough-proof},
\ref{sec:soft} and~\ref{sec:full}, following those proofs with their page references.
The comparison of [O, Lemma~3.1 and Corollaries~3.2--3.3, pp.~10--12] is replaced by
Section~\ref{sec:comparison}. The classical inputs are the Matom\"aki--Radziwi\l\l--Tao short exponential-sum
and distance bounds \cite[Theorem~1.7 and (1.12), p.~6]{MRT}, the fundamental lemma
\cite[Theorem~3.6(a), p.~38]{Ford}, Liouville Bombieri--Vinogradov over all residues
\cite[(1.5), p.~2]{PintzBV}, Shiu's positive mean bound \cite[p.~1]{Henriot}, the Nair--Tenenbaum upper theorem
\cite[Theorem~1, pp.~1--2]{Henriot}, bounded independence \cite{Bazzi,Razborov,Braverman}, Hal\'asz in the form
\cite[Chapter~III.4]{Tenenbaum}, and the fixed-modulus prime number theorem and Mertens' theorem \cite{MV}.
The zero-free half-plane $\Re s>7/8$ is not used in any of the four theorems. The classical zero-free region enters
Theorem~\ref{thm:full}, as well as the other signed graph estimates, through the cited MRT distance bound.

The roughness of the target vertex is put into the padding probability, with row cutoff $KW_1/L$.
Different paddings retain three local densities; equal paddings require a separate diagonal term with exponent
$24/17$. Closed walks are rewritten using arrival indicators before the padding sum. The resulting truncation
event is a depth-three Boolean circuit; bounded independence compares it with the independent residue model.
For soft weights, $\sqrt b$ is placed in the operator and $\sqrt b$ in the test vector.

The shift $h$ is fixed throughout; the statements have no exceptional set of $X$ or of scales.
All absolute constants in the theorems and their thresholds $X_0(h)$
(or $X_0(h,r)$ for the fixed soft weight) are ineffective. No statement here prescribes a fixed number of prime factors
or asserts primality. One remaining two-point estimate toward a fixed number of prime factors is
$|\sum_{n\le X}\lambda(n)\lambda(n+h)J_\sigma(n)J_\sigma(n+h)|\le
(1-2\epsilon)\sum_{n\le X}J_\sigma(n)J_\sigma(n+h)$,
where $J_\sigma(n)=\prod_{p\mid n}(1-p^{-\sigma})$, $\sigma=c/\log X$, $0<\epsilon\le1/4$ and $c$ is fixed
sufficiently large in terms of $\epsilon$.

Section~\ref{sec:setup} specifies the inputs and graph data. Section~\ref{sec:comparison} proves finite comparison.
Sections~\ref{sec:collision}--\ref{sec:rough-proof} prove Theorem~\ref{thm:rough}.
Section~\ref{sec:soft} treats the balanced soft weight. Section~\ref{sec:consumer} proves
Theorem~\ref{thm:omega-consumer} and Remark~\ref{rem:tradeoff}.
Sections~\ref{sec:full}--\ref{sec:full-proof} prove Theorem~\ref{thm:full}.

\section{Notation, classical inputs and graph data}\label{sec:setup}
Write $\e(\theta)=\exp(2\pi i\theta)$, $\tau_k$ for the $k$-fold divisor function, $\tau=\tau_2$,
and $P(z)=\prod_{p\le z}p$. Let $\mu$ be the M\"obius function, $\varphi$ Euler's totient, and $\zeta$ the Riemann zeta function.
Write $P^+(n)$ for the largest prime factor, with $P^+(1)=1$. Powers with exponent zero mean one, also at $r=0$.
In an independent residue model there is one uniform coordinate $R_s$ modulo each $s$ in a pairwise coprime family.
Write $\Prob_*$ and $\E_*$ for probability and expectation under this product law, and $\Prob_I$ for the law induced
by a uniformly chosen integer in a finite interval $I$ of consecutive integers. Every translated site uses the same
coordinates with the indicated offsets. Translated sites are not independent.
Write $\rho(B)$ for the spectral radius of a finite matrix $B$; the Hilbert--Schmidt norm is unnormalized.
For a finite prime set $\mathcal T$, put $\omega_{\mathcal T}(n)=\sum_{p\in\mathcal T}\ind_{p\mid n}$.
The same divisibility notation is used for a residue site. At integer or residue sites, the rough mask
$R_z(n)$ means $\prod_{p\le z}\ind_{p\nmid n}$.

\begin{lemma}[Rough densities and exceptional primes]\label{lem:density}
Let $W_1(z)=\prod_{p\le z}(1-1/p)$ and $v=\log(2z)$. Then
\begin{equation}\label{eq:rough-density}
\begin{aligned}
 W_1(z)&\asymp v^{-1},\qquad V_h(z)\asymp_h W_1(z)^2,\\
 V_h(z)&\ge c_*W_1(z)^2,\qquad W_1(z)/V_h(z)\ll v,\\
 c_*&=\prod_{p>2}\left(1-\frac1{(p-1)^2}\right)>0.
\end{aligned}
\end{equation}
For a nonzero integer $\Delta$,
\begin{equation}\label{eq:exceptional-harmonic}
 \sum_{p\mid\Delta}\frac1p\le\log\log\log(3|\Delta|)+O(1),\qquad
 \prod_{\substack{p>3\\p\mid\Delta}}(1-3/p)^{-1}
 \ll(\log\log(3|\Delta|))^6.
\end{equation}
\end{lemma}
\begin{proof}
Mertens' theorem \cite[Section~2.2]{MV} gives the assertion about $W_1$.
For $p\nmid h$, $(1-2/p)/(1-1/p)^2=1-1/(p-1)^2$; for $p\mid h$ the ratio is
$(1-1/p)^{-1}$. Since $2\mid h$, the factor at $2$ is positive and all other exceptional factors form a fixed
finite product. This proves every comparison in \eqref{eq:rough-density}.
For \eqref{eq:exceptional-harmonic}, separate primes at $\log(3|\Delta|)$. The smaller ones contribute at most
$\log\log\log(3|\Delta|)+O(1)$. There are at most
$\log(3|\Delta|)/\log\log(3|\Delta|)$ larger prime divisors, and their reciprocal sum is bounded.
Finally $-\log(1-3/p)\le6/p$ for $p\ge5$. Exponentiation gives the product bound.
\end{proof}

\begin{inputthm}[Fundamental lemma]\label{input:fl}
For a sieve problem with multiplicative densities $0\le g(p)<1$ satisfying the product condition of dimension
at most two, there are upper and lower divisor weights $\xi_d^\pm$, $|\xi_d^\pm|\le1$, supported on
squarefree $d\mid P(z)$, $d\le D$, such that the corresponding divisor sums bound the sifted indicator above
and below and
\begin{equation}\label{eq:fl}
 \left|\sum_d\xi_d^\pm g(d)-\prod_{p\le z}(1-g(p))\right|
 \le E(s)\prod_{p\le z}(1-g(p)),\qquad s=\frac{\log D}{\log z},\qquad E(s)\ll\e^{-cs}
\end{equation}
for sufficiently large $s$, when $2\le z\le D^{1/2}$.
The constants are uniform for subsets of the forbidden roots $0,-h$.
This is \cite[Theorem~3.6(a), p.~38; probability-space formulation, pp.~1--3]{Ford}; its sharper error is
$\exp(-s\log s+s\log_3s+O(s))$. Indeed $g(p)\le\nu_h(p)/p$, and the associated product condition is bounded by
the dimension-two product. In particular its constants can be absolute for even $h$.
\end{inputthm}

\begin{lemma}[Maximal all-residue Liouville bound]\label{lem:bv}
For fixed $\vartheta<1/2$ and every $A>0$,
\begin{equation}\label{eq:bv}
 \sum_{d\le T^\vartheta}\max_{a\bmod d}\max_{y\le T}
 \left|\sum_{\substack{n\le y\\n\equiv a\pmod d}}\lambda(n)\right|
 \ll_{A,\vartheta} T(\log T)^{-A}.
\end{equation}
The maximum includes nonunit classes. For $D\le X^{1/4}$, with $E_d$ denoting this maximum at $T=2X$,
\begin{equation}\label{eq:root-cauchy}
 \sum_{d\le D}\mu^2(d)\nu_h(d)E_d
 \le\left(\sum_{d\le D}E_d\right)^{1/2}
     \left(\sum_{d\le D}\mu^2(d)4^{\omega(d)}E_d\right)^{1/2}
 \ll_{h,B}X(\log X)^{-B}
\end{equation}
for every $B>0$.
\end{lemma}
\begin{proof}
The bound without the maximum over $y$ is \cite[(1.5), p.~2]{PintzBV}.
Fix $\vartheta<\vartheta'<1/2$ and an integer $m>A+3$. Use $O(\log^m T)$ cutoffs with spacing at most
$T/\log^m T$. At grid points $y\ge T/(2\log^m T)$ one has $T^\vartheta\le y^{\vartheta'}$ and
$\log y\asymp\log T$. The cited estimate with saving $A+m+3$, summed over the grid, is
$O(T\log^{-A-2}T)$. Interpolation and the initial small interval cost
$O(T\log^{1-m}T+T^\vartheta)$, by the progression term count
$T/(d\log^m T)+1$. This proves \eqref{eq:bv}.
Since $\nu_h(d)\le2^{\omega(d)}$ on squarefree integers, Cauchy--Schwarz gives the inequality in
\eqref{eq:root-cauchy}. Also $E_d\le2X/d+1$ and
\begin{equation}\label{eq:divisor-bounds}
 \sum_{d\le D}\frac{4^{\omega(d)}}d\le(1+\log D)^4,\qquad
 \sum_{d\le D}4^{\omega(d)}\ll D(1+\log D)^3.
\end{equation}
These follow from $4^{\omega(d)}\le\tau_4(d)$ by summing its four divisor variables.
The second square in \eqref{eq:root-cauchy} is $O(X\log^4(2X))$; arbitrary $A$ in \eqref{eq:bv} finishes the proof.
Replacing $n$ by $n+h$ costs two prefixes, once $X\ge h$.
\end{proof}

\begin{lemma}[Unsigned and singly signed rough sums]\label{lem:rough-sieve}
For each fixed $0<\rho_*<1$, uniformly for $2\le z\le\exp((\log X)^{\rho_*})$,
\begin{align}
 \sum_{n\le X}R_z(n)R_z(n+h)&=XV_h(z)+O_{h,\rho_*,B}(X\log^{-B}X),\label{eq:rough-unsigned}\\
 \sum_{n\le X}\lambda(n)R_z(n)R_z(n+h),\quad
 \sum_{n\le X}\lambda(n+h)R_z(n)R_z(n+h)&\ll_{h,\rho_*,B}X\log^{-B}X.\label{eq:rough-single}
\end{align}
Both assertions hold for every fixed $B>0$.
\end{lemma}
\begin{proof}
Put $D=X^{1/4}$, $g_h(d)=\nu_h(d)/d$ on squarefree $d$, and
$w^\pm(n)=\sum_{d\mid n(n+h)}\xi_d^\pm$. CRT gives
\begin{equation}\label{eq:sieve-crt}
 \sum_{n\le X}\ind_{d\mid n(n+h)}=Xg_h(d)+O(\nu_h(d)).
\end{equation}
Since $\sum_{d\le D}\mu^2(d)\nu_h(d)\ll D\log(2D)$, Input~\ref{input:fl} yields
\begin{align}
 \sum_{n\le X}R_z(n)R_z(n+h)&=XV_h(z)+O(E(s)XV_h(z)+D\log(2D)),\label{eq:sieve-count}\\
 \sum_{n\le X}|w^+(n)-R_z(n)R_z(n+h)|&\le2E(s)XV_h(z)+O(D\log(2D)),\label{eq:sieve-lone}\\
 \qquad s&\ge\tfrac14(\log X)^{1-\rho_*}.\notag
\end{align}
For large $X$, $z\le D^{1/2}$, $g_h(p)<1$ because $h$ is even, and $E(s)$ beats every logarithmic power.
The second bound follows from $w^+\ge R_zR_z\ge w^-$. It is an $L^1$ bound and may be multiplied by either
Liouville sign. Expanding $w^+$ and applying \eqref{eq:root-cauchy} proves the two signed bounds.
\end{proof}

\begin{inputthm}[Short Fourier means and distance]\label{input:mrt}
For a multiplicative $g$ with $|g|\le1$, $10\le D\le Y$, and
$Q=\min((\log Y)^{1/125},(\log D)^5)$, define
\begin{equation}\label{eq:pretentious}
 M(g;Y,Q)=\min_{\substack{q\le Q,\ \chi\bmod q\\|t|\le Y}}
 \sum_{p\le Y}\frac{1-\Re(g(p)\overline{\chi(p)}p^{-it})}p.
\end{equation}
Then
\begin{equation}\label{eq:mrt}
 \sup_\theta\frac1Y\int_0^Y
 \left|\sum_{x\le n\le x+D}g(n)\e(\theta n)\right|dx
 \ll D\left(\e^{-M(g;Y,Q)/20}+\frac{\log\log D}{\log D}+(\log Y)^{-1/700}\right).
\end{equation}
Moreover in this character and height range
\begin{equation}\label{eq:lambda-distance}
 M(\lambda;Y,Q)\ge\tfrac14\log\log Y-O(1).
\end{equation}
These are \cite[Theorem~1.7 and (1.12), p.~6]{MRT}, with a conservative choice of the constant in (1.12).
For the distance assertion one can apply (1.12) with auxiliary long and short parameters both $Y$ and its
coefficient bound $1$, giving the distance at prime cutoff $10Y$ for every $q\le Q$ and $|t|\le Y$.
Removing the primes in $(Y,10Y]$ costs at most $2\sum_{Y<p\le10Y}1/p=O(1)$, giving \eqref{eq:lambda-distance}.
The supremum in \eqref{eq:mrt} is outside the origin integral. The distance bound in the cited version is deduced
from the classical Vinogradov--Korobov zero-free region.
\end{inputthm}

\begin{definition}[Labels, bins and vertex cuts]\label{def:graph}
Fix $W\ge10$, $J=\lfloor\alpha\log L\rfloor\ge1$, with $\alpha=1/(1200W)$, and $0<\eta\le1$.
Let $H_0=\exp(L^{199/200})$, $Y_i=L^{199/200}\e^{6Wi}$, and
\begin{equation}\label{eq:prime-supplies}
 \Pset_i=\{p\equiv1\pmod5:Y_{i-1}\le\log p<Y_i,\ p\nmid h\},\quad
 \Pset=\bigcup_{i=1}^J\Pset_i,\quad
 \Dset=\{p_1\cdots p_J:p_i\in\Pset_i\}.
\end{equation}
At a small-prime cutoff $z$ (or $y$) below $H_0$ take
\begin{equation}\label{eq:padding-data}
\begin{gathered}
 \Qset=\{p:z<p\le\e^L,\ p\not\equiv1\pmod5,\ p\nmid h\},\\
 a(u)=a^{\omega(u)},\quad w(n)=(1+a)^{\omega_{\Qset}(n)},\quad S=\prod_{p\in\Qset}(1+a/p).
\end{gathered}
\end{equation}
where $a=4$ or $16$, and $\Qset^{\sfw}$ is the finite set of squarefree products of primes in $\Qset$, including $1$.
Set
\begin{equation}\label{eq:bins}
 \mathcal S_j=\{(d,u)\in\Dset\times\Qset^{\sfw}:\omega(u)\le100\log L,
 j\eta\le\log(du)<(j+1)\eta\},\quad T_j=X\e^{j\eta},\quad j\eta\le100L.
\end{equation}
Write $V_i=\sum_{p\in\Pset_i}1/p$, $V_*=\prod_iV_i$,
$S_0=\sum_j\sum_{(d,u)\in\mathcal S_j}a(u)/(du)$, and $k=\lfloor L\rfloor$,
$\ell_0=\lfloor L^{1/10}\rfloor$, $M=\lceil\e^{103L}\rceil$.
The exterior projection $E$ restricts to $\omega_{\Pset}(n)\le6WJ$.
\end{definition}

\begin{proposition}[Label supply and retained mass]\label{lem:mass}
With the fixed data above, for sufficiently large $L$,
\begin{equation}\label{eq:label-mass}
 W\le V_i\le2W,\quad \log d<2L,\quad
 \tfrac12SV_*\le S_0\le SV_*,\quad Xe^{-\eta}<T_j/(du)\le X.
\end{equation}
\end{proposition}
\begin{proof}
Since $6WJ\le(1/200)\log L$, the supplies lie below $\e^L$. The fixed-modulus prime number theorem
\cite[(2.2)--(2.3), pp.~8--9]{OpenAI007} gives
$V_i=(3/2)W+o(1)$ uniformly in $i$; removing the fixed divisors of $h$ has bounded, eventually vanishing effect.
The geometric sum of $Y_i$ is less than $L/(1-e^{-6W})<2L$.
Under the law $a(u)/(uS)$, a prime occurs with probability $a/(p+a)$.
For $a=4$ the expected logarithmic size and degree are at most $4L$ and $4\log L$.
For $a=16$ the bounds $17L$ and $17\log L$ suffice for large $L$, also after removing primes up to $z$.
Markov excludes $\log u>98L$ and $\omega(u)>100\log L$ with probabilities totaling at most
$4/98+4/100<1/2$ or $17/98+17/100<1/2$, respectively.
Every remaining $u$ belongs to a bin with every $d$. This repeats the two exclusions in
\cite[Lemma~2.2, pp.~8--9]{OpenAI007}, with the present harmonic law proved explicitly.
Dropping restrictions gives the upper bound. The last assertion follows directly from \eqref{eq:bins}.
\end{proof}

\begin{definition}[Rare vertices]\label{def:rare}
A signed numerical step is $D_i=\epsilon_i h d_i u_i$, $\epsilon_i\in\{\pm1\}$, $(d_i,u_i)\in\mathcal S_j$.
A path is positive when $d_i u_i$ divides its departure at each step; this also holds at the arrival.
A forward prohibited word has length $3\le m\le\ell_0$, unequal consecutive whole tuples,
interval-shaped occurrences of every centre prime, and a prime $p\mid d_1$, $p\nmid d_m$, such that
\begin{equation}\label{eq:prohibited}
 p\mid\sum_{i=b}^m D_i\qquad\text{for some }1<b<m.
\end{equation}
It is minimal if no shorter contiguous subword is forward prohibited in either orientation.
Let $\mathcal B_j$ be the sites where a positive minimal prohibited word starts, and $\sigma_j=\ind_{n\notin\mathcal B_j}$.
These are the definitions of \cite[Definition~5.1, pp.~21--22]{OpenAI007} with the present label supply.
\end{definition}

\begin{proposition}[Rare-site estimate for the changed padding]\label{lem:rare}
For either padding weight in \eqref{eq:padding-data},
\begin{equation}\label{eq:rare}
 \Prob_*(n\in\mathcal B_j)\le\exp(-\tfrac12L^{199/200}).
\end{equation}
Its positivity event is a disjunction of at most $\exp(O(\ell_0 L))$ conjunctions, each with
$O(\ell_0\log L)$ residue equalities.
\end{proposition}
\begin{proof}
Follow \cite[Lemmas~5.2--5.3, pp.~22--23]{OpenAI007}. A word of at most $R\le4L$ steps has
$O(R\log L)$ prime slots, so its equality patterns and polynomially many records per slot number
$\exp(O(R\log^2L))$. Each distinct prime contributes one reciprocal and harmonic mass $O(\log L)$.
Now $La(u)\le L^{1+100\log a}$; this changes only the absolute constant in that count.
In a prohibited word, a centre prime $z'$ entering on the last transition has coefficient
$\epsilon_m h u_m d_m/z'$ in \eqref{eq:prohibited}. The controlling prime $p$ is absent from $d_m$, divides
neither $h$ nor any padding, and so that coefficient is a unit modulo $p$, independent of $z'$.
With the other labels fixed the reciprocal sum for $z'$ is bounded by
\begin{equation}\label{eq:elimination}
 \sum_{\substack{H_0\le n\le e^L\\n\equiv c\pmod p}}\frac1n
 \le\frac1{H_0}+\frac Lp\le\frac{1+L}{H_0}.
\end{equation}
The remaining positive reciprocal sum is $\exp(O(\ell_0\log^2L))$.
Thus the probability is at most $(1+L)H_0^{-1}\exp(O(\ell_0\log^2L))$, giving \eqref{eq:rare}.
Numerical pairs have $du\le e^{101L}$ and are uniquely recovered from their product by disjoint supports.
Preselect the numerical prohibited and minimal words; positivity alone is a conjunction of the stated size.
This proves the Boolean description and translation covariance. No Liouville values enter this event.
\end{proof}

\section{Finite residue comparison}\label{sec:comparison}

\begin{lemma}[Shallow CNF comparison]\label{lem:cnf}
Fix an integer $\ell\ge1$. Let the coordinate moduli be pairwise coprime integers between $2$ and $e^L$,
all coprime to $\ell$, with the coordinate modulo $\ell$ added if $\ell>1$.
An event given by a residue CNF with at most $e^{L^3}$ clauses, each of width at most $L$, has
\begin{equation}\label{eq:cnf-comparison}
 |\Prob_I(E)-\Prob_*(E)|\le e^{-L^6}\qquad(|I|\ge e^{L^{18}/2}).
\end{equation}
The literals may be residue equalities or their negations, with repetitions and arbitrary offsets.
An expansion of such events with total absolute coefficient mass $e^{CL^4}$ has error at most $e^{-L^5}$.
The threshold depends only on $\ell,C$.
\end{lemma}
\begin{proof}
Bazzi's bounded-independence theorem \cite{Bazzi}, in the form of \cite[Theorem~1, p.~2]{Razborov}, bounds the
error for an $m$-clause CNF under exactly $t$-wise uniform bits by
$m^{C_0}e^{-c_0\sqrt t}$; complementation gives CNFs from DNFs.
We give the reduction to that precise hypothesis.

Put $B_0=\lceil L^8\rceil$. Adjoin an independent $U_s$ uniform on $[0,1)$ at every coordinate, and encode
$K_s=\lfloor2^{B_0}(R_s+U_s)/s\rfloor$ in $B_0$ bits. Decode by
\begin{equation}\label{eq:decoder}
 \widehat R_s=\left\lfloor\frac{s(K_s+1/2)}{2^{B_0}}\right\rfloor,\qquad
 \Prob(\exists s:\widehat R_s\ne R_s)
 \le2^{1-B_0}\left(\ell+\sum_s s\right)\le e^{3L}2^{-B_0}.
\end{equation}
Conditional on any residue, only its two boundary dyadic cells can cause failure, of probability at most
$2s2^{-B_0}$. Under the product law all encoded bits are independent and unbiased.
The monotone decoder makes every residue preimage an interval of binary integers.
An interval has a CNF with at most $2B_0$ clauses, by excluding each forbidden initial differing bit at each
boundary. Its complement has at most $B_0^2$ clauses after distributing the two one-sided CNFs.
Thus a residue literal has at most $2B_0^2$ clauses. Distributing a clause of width at most $L$ gives a Boolean CNF
with
\begin{equation}\label{eq:cnf-size}
 m\le e^{L^3}(2B_0^2)^L\le e^{L^4},\qquad n_b\le e^{2L}
\end{equation}
clauses and bits, respectively, for large $L$.

Take $t=\lceil L^{13}\rceil$. At most $t$ selected bits read at most $t$ coordinates, whose product is at most
$\ell e^{tL}$. CRT counts each residue with error $O(1)$, so their total variation from uniform is at most that
product divided by $|I|$, and hence at most $\Delta=e^{2tL-L^{18}/2}$.
Jitter and projection do not increase this distance. If $f$ is the full bit-law density relative to the uniform
cube, its Walsh coefficients satisfy $|\widehat f(S)|\le2\Delta$ for $1\le|S|\le t$. Consequently
\begin{equation}\label{eq:walsh-mass}
 q:=\sum_{1\le|S|\le t}|\widehat f(S)|\le2\Delta(t+1)n_b^t\le e^{-L^{17}}.
\end{equation}
Let $H=\sum_{1\le|S|\le t}\widehat f(S)\chi_S$ and $a_0=e^{-2L^6}$. Since $|H|\le q\le a_0$,
\begin{equation}\label{eq:walsh-correction}
 g=\frac{f-H+a_0}{1+a_0}\ge0,\qquad \E g=1,\qquad
 d_{\rm TV}(f,g)\le\frac{q+a_0\E|1-f|}{2(1+a_0)}\le3a_0/2.
\end{equation}
All nonconstant Walsh coefficients of order at most $t$ vanish for $g$, so Fourier inversion on each marginal
makes $g$ exactly $t$-wise uniform.
Bazzi's error is at most $\exp(C_0L^4-c_0L^{13/2})\le a_0$.
Adding the correction and decoding errors gives
$2e^{3L}2^{-B_0}+5a_0/2\le e^{-L^6}$.
Multiplication by $e^{CL^4}$ gives the scalar assertion. The argument used only the coordinate count, sizes and
CRT, so it applies to prime-square coordinates as well as primes.
\end{proof}

\begin{proposition}[Encoding the graph events]\label{prop:encoding}
For padding $4^{\omega(u)}$, any padding-prime subset as above, and any fixed-power choices of $\eta,K$, every
comparison used in positive deletions, the signed trace, and the truncated testing norm is a scalar expansion
covered by Lemma~\ref{lem:cnf}, of total absolute coefficient mass $\exp(O(L^4))$.
Small-prime rough indicators or prime-square negative tests may be adjoined.
\end{proposition}
\begin{proof}
This replaces \cite[Lemma~3.1 and Corollaries~3.2--3.3, pp.~10--12]{OpenAI007} for the present objects.
Follow the event expansions of \cite[Lemma~4.6, pp.~19--21; Theorem~5.4, pp.~24--25;
Proposition~8.4, pp.~41--42]{OpenAI007}.
A numerical witness is preselected by its numerical arithmetic and minimality conditions.
Its positivity is a conjunction of $O(L^{1/10}\log L)<L$ equalities. There are
$\exp(O(L^{11/10}))$ witnesses at one site; their absence is a CNF. Conjoining these tests at at most
$2k\le2L$ visited sites still gives fewer than $e^{L^3}$ clauses.

An exact active-$\Qset$ state is a conjunction of all presences and all absences, with unit-width clauses.
After the cutoff $\omega_{\Qset}\le\lfloor400\log L\rfloor$, there are
$\exp(O(L\log L))$ possible states at a site. In each, denominators and the original padding cutoff are scalars.
A degree condition $\omega_{\mathcal T}\le b$, $b=O(\log L)$, excludes each simultaneous $(b+1)$-subset and is a
CNF of $\exp(O(L\log L))$ clauses. Roughness and fixed divisibilities add unit clauses.

In a trace, fix a signed numerical word and the two external copy indices. Enumerate its states and expand its
$2kJ$ centred factors into indicators and scalars, retaining their signs. The resulting event is the conjunction
just described. There are $\exp(O(kL\log L))$ states, at most $2^{2kJ}$ centred expansions, scalar sizes
$\exp(O(k\log L))$, and $\exp(O(kL))$ numerical words and external indices.
Their total mass is $\exp(O(L^2\log L))\le\exp(O(L^4))$.
This comparison precedes absolute values of word expectations.

For deletion, expand the positive centred majorant and split the four failed cuts at both endpoints by the
exact union bound. A large-degree or witness event $B$ may be treated by
\begin{equation}\label{eq:two-cnf}
 \ind_{H\cap B}=\ind_H-\ind_{H\cap B^c}.
\end{equation}
Both right-hand events have the required CNFs. This avoids an expansion over intersections of arbitrary witness
families. Eligible pairs number $e^{O(L)}$ and exact states cost $e^{O(L\log L)}$, within the same scalar bound.
Finally the truncated norm is expanded by its exact active state, of coefficient $5^{|\Qset_n|}$, at most a fixed
power of $L$. Its $M\le e^{104L}$ block positions preserve the bound.
Apply Lemma~\ref{lem:cnf} to all expansions. None reads the full $\Omega$ weight or a Liouville value.
\end{proof}

\begin{proposition}[Rough-target threshold comparison]\label{prop:threshold}
For padding $16^{\omega(u)}$ and $w=17^{\omega_{\Qset}}$, replace the vertex degree cutoff by
$\lfloor18\log L\rfloor$. The norm, positive deletion and signed-trace comparisons with the rough-target
threshold defined in \eqref{eq:masked-rho} have per-event error $e^{-L^{10}}$ and aggregate error $e^{-L^9}$ on
any origin interval of length at least $e^{L^{625}/2}$. These estimates are uniform in
$z\le e^{L^{1/6}}$, bin position and numerical word data, for fixed $h$.
\end{proposition}
\begin{proof}
On an exact active state there are at most
$N=2^{18\log L}=L^{18\log2}<L^{13}$ possible padding divisors.
A weighted sum of their target rough indicators exceeds its threshold exactly when some subset of candidates
whose scalar weights exceed the threshold has all its targets rough. Negating this statement gives an AND over
those subsets of an OR of positive residue equalities: one selected target has a small prime divisor.
There are at most $2^N$ clauses, each of width at most $N\pi(z)$. Both orientations, all $O(L)$ visited sites,
exact states, no-witness predicates and endpoint masks give a residue CNF whose logarithmic size is
$O(L^{18\log2})+O(L^{1/6})+O(L\log L)$.
Use \eqref{eq:two-cnf} for positive bad-event costs. The subsets are internal to this circuit, rather than scalar
summands.

Use the monotone decoder \eqref{eq:decoder} with $B_0=\lceil L^{15}\rceil$.
Each binary interval has a DNF with at most $2B_0$ dyadic-prefix conjunctions, by following its two boundary
paths in the binary tree and taking the included subtrees. Its complement has $O(B_0)$ such conjunctions too.
Substitution into the residue CNF and flattening adjacent OR levels gives AND--OR--AND, of depth three and size
at most $e^{L^{14}}$, including input occurrences. Decoding error is charged per coordinate: one correct decode
settles every offset test at that coordinate.

Braverman's exact independence bound \cite[Corollary~2, p.~1]{Braverman} for depth $d$, size $m$ and error
$\epsilon_0$ is
\begin{equation}\label{eq:braverman}
 3\,60^{d+3}(\log m)^{(d+1)(d+3)}
       \{5\log(15m/\epsilon_0)\}^{d(d+3)}.
\end{equation}
At $d=3$, $m\le e^{L^{14}}$ and $\epsilon_0=e^{-L^{11}}$, this is $O(L^{588})$.
Equivalently the Main Theorem there may be applied with
\begin{equation}\label{eq:braverman-ceiling}
 s_B=\lceil100L^{14}\rceil,\qquad 0.82^{s_B}(15m)\le e^{-L^{11}}
\end{equation}
for large $L$; the same exponent follows independently of the logarithm-base convention.
Take $t=\lceil L^{589}\rceil$. At most $t$ bits involve moduli of product $e^{O(tL)}$, and the total bit count is
at most $e^{2L}$. CRT and the low-order Walsh coefficient count therefore cost
$\exp(-L^{625}/2+O(L^{590}))\le e^{-L^{12}}$.
Apply the nonnegative correction \eqref{eq:walsh-correction} with $a_0=e^{-L^{11}}$ to obtain exact $t$-wise
uniformity. Its error is $O(e^{-L^{11}})$; adding Braverman's error and the coordinate decoding error
$e^{3L}2^{-L^{15}}$ gives $e^{-L^{10}}$.
The states, words and factors $L16^{\omega(u)}$ retain the scalar mass $e^{O(L^4)}$ proved above.
Thus their aggregate error is $e^{-L^9}$. The large threshold is not covered merely by Lemma~\ref{lem:cnf}.
\end{proof}

\section{Rough-target padding and its second moment}\label{sec:collision}
Use $a=16$ in Definition~\ref{def:graph}, and suppose $2\le z\le e^{L^{1/6}}$.
For a numerical $d$ whose primes exceed $z$ and are outside $\Qset$, $\log d\le2L$, define
\begin{equation}\label{eq:masked-rho}
 \rho_{j,z}^{\pm}(n,d)=\frac1{w(n)}
 \sum_{\substack{u\mid n,\ u\in\Qset^{\sfw}\\j\eta\le\log(du)<(j+1)\eta}}
 16^{\omega(u)}R_z(n\pm hdu),\qquad
 \rho_{j,z}=\max(\rho_{j,z}^+,\rho_{j,z}^-).
\end{equation}
The sum here includes all padding divisors in the bin, before the eligible degree restriction.

\begin{proposition}[Three-density second moment]\label{prop:collision}
For $0<\eta\le1$ and the bins $j\eta\le100L$, uniformly in $z,\eta,d$,
\begin{equation}\label{eq:masked-second}
 \E_*\left[w(n)R_z(n)\sum_j\rho_{j,z}(n,d)^2\right]
 \ll_h\frac{S W_1^3(1+\log L)^6}{L}.
\end{equation}
Deleting a row with $\rho_{j,z}>KW_1/L$, after the low-$\Qset$ cutoff, costs, after summing bins and labels and
normalizing by $SV_*$, at most
\begin{equation}\label{eq:masked-deletion}
 C_h W_1^2\frac{2^J(1+\log L)^6}{K}.
\end{equation}
\end{proposition}
\begin{proof}
Following the conditional-divisor experiment of \cite[Lemma~4.5, pp.~17--18]{OpenAI007}, tilt by $w/S$.
Each $\Qset$ prime is then active with probability $17/(p+16)$, and conditional on activity is selected in a
divisor with probability $16/17$. Two conditionally independent selected divisors $u_1,u_2$ give the factor
\begin{equation}\label{eq:collision-factor}
 1-\frac{32}{17(p+16)}(1-\cos(t\log p))
\end{equation}
for their log-ratio characteristic function. It is nonnegative for every real $t$; the harmonic estimate below is used only for $|t|\le2$.
Fixed-modulus PNT and partial summation, as in \cite[(4.11)--(4.13), pp.~17--18]{OpenAI007}, give
\begin{equation}\label{eq:collision-char}
 0\le\phi_z(t)\ll_h
 \left(\frac{1+v|t|}{1+L|t|}\right)^{24/17},\qquad v=\log(2z)\le2L^{1/6},\quad |t|\le2.
\end{equation}
Indeed the harmonic sum of $1-\cos(t\log p)$ on these padding primes is
$\tfrac34\log((1+L|t|)/(1+v|t|))+O_h(1)$; replacing $p+16$ by $p$ changes it by $O(1)$.
The coefficient $(32/17)(3/4)$ is $24/17$.
The nonnegative triangular Fourier kernel
\begin{equation}\label{eq:triangular}
 \kappa(x)=\left(\frac{\sin x}{x}\right)^2
 =\tfrac14\int_{-2}^2(2-|t|)e^{itx}\,dt,\qquad
 \ind_{[-1,1]}(x)\le\kappa(x)/(\sin1)^2
\end{equation}
bounds bin collision. If $\rho_j$ denotes the unmasked bin probability, it gives
\begin{equation}\label{eq:unmasked-collision}
 \E^{\rm tilt}\sum_j\rho_j^2\ll_h L^{-1}+(v/L)^{24/17}\ll_h L^{-1},\qquad
 \E^{\rm tilt}F=\E_*(wF)/S.
\end{equation}
The exact equality probability is
\begin{equation}\label{eq:diagonal-padding}
 D_Q=\prod_{p\in\Qset}\left(1-\frac{32}{17(p+16)}\right)\ll_h(v/L)^{24/17}.
\end{equation}

Fix an orientation and condition on the active padding configuration and $u_1,u_2$.
For equal paddings, the two small-prime sites have mean $V_h(z)\ll_h W_1^2$, since $du$ is a unit at every small
prime. For different eligible paddings, $0<|u_1-u_2|\le e^{101L}$.
At each $p>3$ not dividing $h(u_1-u_2)$, the three offsets $0,\pm hdu_1,\pm hdu_2$ are distinct modulo $p$.
Dropping the other factors bounds their common roughness probability by
\begin{equation}\label{eq:triple-rough}
 \prod_{3<p\le z}(1-3/p)
 \prod_{\substack{p>3\\p\mid h(u_1-u_2)}}(1-3/p)^{-1}
 \ll_h W_1^3(1+\log L)^6.
\end{equation}
Use Lemma~\ref{lem:density}; if the forbidden roots cover a whole small field, the probability is zero.
Apply this bound before enlarging the unmasked collision sum, so the difference bound remains valid.
Enlarge only the diagonal to all equal divisor choices. Since $\max(x,y)^2\le x^2+y^2$, the left side of
\eqref{eq:masked-second} is bounded by a constant times
\begin{equation}\label{eq:second-diagonal}
 S\left\{\frac{W_1^3(1+\log L)^6}{L}+W_1^2(v/L)^{24/17}\right\},\qquad
 \frac{W_1^2(v/L)^{24/17}}{W_1^3/L}
 \ll_h v^{41/17}L^{-7/17}\ll_h L^{-1/102}.
\end{equation}
The last exponent is $(1/6)(41/17)-7/17=-1/102$. This proves \eqref{eq:masked-second}.

On a bad row, the positive source edge mass $wR_z\rho_{j,z}^+$ is at most
$LwR_z\rho_{j,z}^2/(KW_1)$. The centre majorant
$a_d^+(n)=\prod_{p\mid d}(\ind_{p\mid n}+1/p)$ has mean $2^J/d$ independently of these coordinates.
Sum $d$ to obtain $2^JV_*$ and use \eqref{eq:masked-second}. Reversing the orientation proves the arrival bound.
This gives \eqref{eq:masked-deletion}, without independence between translated sites.
\end{proof}

\section{The modified graph moment and spectral estimate}\label{sec:moment}
For the rough operator let
\begin{equation}\label{eq:rough-tau}
 \tau_j(n,d)=\ind_{\omega_{\Qset}(n)\le\lfloor18\log L\rfloor}
 \ind_{\rho_{j,z}(n,d)\le KW_1/L}\sigma_j(n),\qquad
 c_d(n)=\prod_{p\mid d}(\ind_{p\mid n}-1/p),\quad a_d(n)=|c_d(n)|.
\end{equation}
On $M$ consecutive sites define a symmetric matrix, for $m=n+hdu$ with $(d,u)\in\mathcal S_j$, by
\begin{equation}\label{eq:rough-matrix}
 A_d(n,m)=\frac{L16^{\omega(u)}}{\sqrt{w(n)w(m)}}\ind_{u\mid n}c_d(n)
 \tau_j(n,d)\tau_j(m,d)R_z(n)R_z(m),\qquad A_d(m,n)=A_d(n,m).
\end{equation}
Other entries vanish. On the direct sum of $d$-indexed copies set
$H_{d,d'}=\ind_{d\ne d'}A_{d'}$. The $\Pset$-degree restriction occurs only in $E$.


For endpoint functions $f,g$ and a numerical sign $\epsilon(du)$ define the centred sums
\begin{equation}\label{eq:centered-sums}
 C_j(f,g)=\frac1{T_j}\sum_{n\le T_j}\sum_{(d,u)\in\mathcal S_j}
 a(u)\ind_{u\mid n}c_d(n)\epsilon(du)f(n)g(n+hdu).
\end{equation}
Their retained versions $C_j^\circ(f,g)$ insert $\chi_j(n,d)\chi_j(n+hdu,d)$, where
$\chi_j=\tau_j\ind_{\omega_{\Pset}\le6WJ}$.
For the rough case take $f=g=\lambda R_z$, $a=16$ and $\epsilon=1$.
For the soft and full weights the endpoint functions, numerical sign, and vertex cuts are specified below.

\begin{inputthm}[Finite-dimensional transfer, {\cite[Lemma~8.1, p.~39]{OpenAI007}}]\label{input:transfer}
Let $B_i$ be self-adjoint operators on a finite-dimensional complex Hilbert space, $P$ an orthogonal projection,
and $\mathcal H_{i,j}=\ind_{i\ne j}B_j$. If $a,b\ge0$, $\|B_i\|\le a$, and
$\sum_i\|B_iv\|^2\le b^2\|v\|^2$ for $v=Pv$, then
\begin{equation}\label{eq:transfer}
 \left\|P\left(\sum_iB_i\right)P\right\|\le3\max(a,b,\rho(\mathcal H)).
\end{equation}
Neither different $B_i$ nor $P$ need commute.
\end{inputthm}

\begin{proposition}[Moment with the rough density]\label{prop:masked-moment}
There are absolute constants $C_2,C$, fixed before $W$, for padding $16$, with the following property.
For each fixed even $h\ge2$ and fixed $W\ge10$, take $X$ sufficiently large in terms of $h,W$.
Use the operator \eqref{eq:rough-matrix}, $\alpha=1/(1200W)$, and
\begin{equation}\label{eq:masked-graph-parameters}
 L=(\log X)^{1/625},\quad 2\le z\le e^{L^{1/6}},\quad
 J=\lfloor\alpha\log L\rfloor,\quad\eta=L^{-\alpha/100},\quad
 K=L^{\alpha\log2+\alpha/100}.
\end{equation}
Then
\begin{equation}\label{eq:masked-moment}
 \E_*\|H^k\|_{\HS}^2\le\{KW_1(C_2\sqrt W)^J\}^{2k}.
\end{equation}
The same bound holds for finite origin intervals of length at least $e^{L^{625}/2}$, after one absolute
enlargement of $C_2$. Outside a fraction at most $e^{-2k}$ of those origins,
\begin{equation}\label{eq:masked-spectrum}
 \left\|E\left(\sum_dA_d\right)E\right\|\le KW_1(C\sqrt W)^J.
\end{equation}
\end{proposition}
\begin{proof}
We check the proof of \cite[Theorem~5.4, pp.~24--38]{OpenAI007} in its order.
Write $\mathfrak b=KW_1$. Expanding $\|H^k\|_{\HS}^2$ joins two paths with common endpoints and reverses one.
The result is a closed word of $2k$ steps with unequal consecutive whole tuples within each half. There are at
most $|\Dset|^2M\le e^{108L}$ external choices, as in [O, (5.4), p.~24]; there is no additional copy choice per
step. Proposition~\ref{prop:threshold} compares this signed expansion before taking absolute values.

On a closed word the product of the two endpoint indicators is the product of the arrival indicators, counting
visits with multiplicity, because $R_z^2=R_z$. Its square-root denominators multiply to one $w$ at each departure.
Make these identities before dropping closure. The contribution now has the form [O, (5.6)--(5.7), p.~25], with
$R\ge0$ independent of all $\Pset$ residue draws, and $G=\prod_i\sigma_j(n_i)$.
The new masks and cutoff may depend on numerical centre labels, but not on their residue draws; they remain in
$R$. Only $G$ contains the rare-witness predicates. This verifies the decisive hypothesis of
[O, Lemma~5.5, p.~26]. A centre prime appearing in exactly one step is a singleton.
At every nonsingleton occurrence expand the centred factor into its indicator (lit) or its scalar $-1/p$ (unlit).
Lit consistency requires all the lit departure residues of one prime to agree.
Force the common residue for each lit nonsingleton. For each subset $H$ of the singleton set $\mathcal U$, also
force the residue at the sole occurrence of each singleton in $H$, leaving the other coordinates at their original
draws. Write $G_H$ for $G$ on this hybrid and define
\begin{equation}\label{eq:mixed-difference}
 \Delta G=\sum_{H\subset\mathcal U}(-1)^{S_1-|H|}G_H,\qquad S_1=|\mathcal U|.
\end{equation}
In particular, at a singleton prime,
\begin{equation}\label{eq:singleton-identity}
 \E[(\ind_{Z=a}-1/p)F(Z)]=p^{-1}\E[F(a)-F(Z)].
\end{equation}
Integrating all singleton coordinates gives the same mixed difference $\Delta G$, bounded by $2^{S_1}$,
and the reciprocal factor
\begin{equation}\label{eq:lit-reciprocals}
 b_{\mathcal L}=\prod_{p\ \mathrm{singleton}}p^{-1}
 \prod_{p\ \mathrm{nonsingleton}}p^{-u_p-\ind_{\ell_p>0}},\qquad
 |\E(\text{designated term})|\le\ind_{\mathrm{lit\ consistency}}b_{\mathcal L}\E[R|\Delta G|].
\end{equation}
Write $\mathcal L$ for the lit/unlit designation. Here $\ell_p,u_p$ count lit and unlit occurrences,
and $U=\sum_{p\ \mathrm{nonsingleton}}u_p$ is the total number of unlit nonsingleton occurrences.
Lit tests must force the same residue at their departures.
This is a direct application of the numbered source lemma, with its independence condition verified.

Next follow the three discarded classes. Let $\mathcal M_a$ be the crude factor in the proof of
[O, Lemma~5.2, p.~22] and the discussion preceding Lemma~5.6 (p.~26), with padding coefficient $a^{\omega(u)}$.
It retains the padding divisibilities, but drops denominators, cutoffs and added small-prime masks. Pointwise,
for the non-$\Pset$ factor of any fixed numerical word,
\begin{equation}\label{eq:crude-domination}
\begin{aligned}
 0\le R_{\rm mask}
 &\le\mathcal M_{16}:=\prod_{i=1}^{2k}L16^{\omega(u_i)}\ind_{u_i\mid n_i},\\
 \mathcal M_{16}
 &=4^{\sum_i\omega(u_i)}\mathcal M_4
 \le\exp(200k(\log4)\log L)\mathcal M_4
 =\exp(O(L\log L))\mathcal M_4.
\end{aligned}
\end{equation}
Here $w\ge1$ and $\omega(u_i)\le100\log L$. The remaining lit-consistency and reciprocal factors are identical.
Thus replacing $4$ by $16$ introduces only this additional factor in the positive majorants of
[O, Lemmas~5.6, 6.2, 6.5, pp.~26, 28--29, 31--33]. Their crude enumeration still costs
$\exp(O(L\log^2L))$, with all padding divisibilities retained. The added small-prime draws integrate to at most
one after their masks are dropped, rather than requiring a sum over their configurations.
For $U>L^{1/50}$ unlit occurrences, at least $U/2$ extra reciprocal powers remain. They give
\begin{equation}\label{eq:unlit-cost}
 \exp(O(L\log^2L))H_0^{-U/2}
 \le\exp(O(L\log^2L)-\tfrac12L^{203/200}),
\end{equation}
as in [O, Lemma~5.6, p.~26].

For high rank use [O, Lemma~6.1, pp.~27--28] as stated. In one column let $z_i$ be the abstract label at step $i$,
with numerical prime value $p_{z_i}$, and let $(\mathbf e_z)$ be the coordinate basis indexed by its abstract labels. Put
\[
 t_i=D_i/p_{z_i},\qquad v(i)=\sum_{a<i}t_a\mathbf e_{z_a}.
\]
For two lit occurrences $i,i'$ of a label $z$, its control/difference pair is
$(\mathbf e_z,v(i)-v(i'))$. Joint independence of a family means independence of all its displayed vectors over $\R$.
The coefficients $t_i$ are independent of the numerical prime values in that column. A jointly independent family of
$r_0=\lfloor L^{1/50}\rfloor$ such pairs has an invertible minor over $\R$ after deleting its control columns.
The independent reciprocal prime law has largest atom at most $(V_iH_0)^{-1}$.
The differences in the minor are bounded by $2kh e^{103L+1}=e^{O(L)}$, since eligible paddings still obey
$u\le e^{100L+1}$. These are precisely the atom, independence and size hypotheses of Lemma~6.1; no invertibility
modulo a control prime is required. Its square-root probability saving is
$\exp(O(r_0\log L))H_0^{-r_0/2}$.
Multiplying the crude enumeration gives $\exp(-cL^{203/200})$.
This redoes the application in [O, Lemma~6.2, pp.~28--29] for the new fixed padding coefficient.

For $S_1>L^{1/4}$ singletons, [O, Lemma~6.3, p.~30] applies: $\Delta G\ne0$ forces
$T=\lfloor L^{1/12}\rfloor$ positive witnesses in a single hybrid, with distinct private marked singletons,
and recorded total length at most $2k+\ell_0T\le4L$.
Its proof uses only the witness supports and the mixed difference, unchanged here.
[O, Lemma~6.4, pp.~31--32] supplies either an active private label with a unit coefficient in an internal
congruence, or an active label with a nonzero private-label prefix contribution. The prohibited-word and
minimality definitions are exactly Definition~\ref{def:rare}, so its statement applies as given.
To repeat [O, Lemma~6.5, pp.~32--33], order witnesses by attachment. If at least $T/8$ have an active label
absent from all earlier witnesses, use their internal congruences and eliminate in reverse order; if absent from
all later witnesses, reverse that order. Otherwise more than $3T/4$ witnesses have all their active labels in both
an earlier and a later witness. For private labels whose main occurrence has not passed their attachment, compare
with an earlier witness; the main segment omits that label, and privacy removes it from the partner witness.
For labels already passed, compare with a later witness and reverse the attachment order.
One group supplies at least $\lfloor T/8\rfloor$ relations. Each relation is independent of all later selected
variables and its selected coefficient is a unit. Use \eqref{eq:elimination} successively. The residue tests for
new witness primes are unchanged in every hybrid, so no hybrid-choice cost is incurred. All records still cost
$\exp(O(L\log^2L))$. The resulting bound is
\begin{equation}\label{eq:singleton-cost}
 \exp(O(L\log^2L))\left(\frac{1+L}{H_0}\right)^{\lfloor T/8\rfloor}
 \le e^{-L^{21/20}},\qquad T\log H_0\asymp L^{647/600}.
\end{equation}
Thus all three exceptional classes have total $\exp(-cL^{1+\epsilon_1})$ for some absolute $\epsilon_1>0$.
The added small-prime draws were dropped in these positive bounds, with no sum over their configurations.

After the three classes are discarded, [O, Lemmas~7.1--7.2, pp.~34--37] apply to the remaining words as stated.
Perfect positions are lit nonsingletons; their blocks have length at most $\ell_0$, have no prohibited witness,
and have unequal consecutive tuples within each half. Their interval-shaped labels and nonzero subinterval
displacements follow from Lemma~7.1. The bounds $S_1\le L^{1/4}$, $U\le L^{1/50}$ and rank below $r_0$ are
exactly the hypotheses of Lemma~7.2. Its forest universe has $e^{O(Jk)}$ combined patterns uniformly over all
numerical labels and padding coefficients. Added masks only select a subset of this universe.

At this point repeat [O, Lemma~7.3 and (7.5), p.~37], changing its row constant.
Fix the non-$\Pset$ environment, numerical tuples $d_i$ and orientations $\epsilon_i\in\{\pm1\}$.
The coefficient at a departure $n$ is
\begin{equation}\label{eq:rough-recursion-coefficients}
\begin{aligned}
 c_i(n,u)={}&\frac{L a(u)}{w(n)}\ind_{u\mid n}\ind_{\rm bin}
 R_z(n+\epsilon_i h d_i u)\\
 &\hspace{1cm}\cdot\ind_{\rho_{j,z}(n,d_i)\le KW_1/L},\qquad a(u)=16^{\omega(u)}.
\end{aligned}
\end{equation}
The retained oriented row satisfies
\begin{equation}\label{eq:density-row}
 \sum_u\frac{L16^{\omega(u)}}{w(n)}\ind_{u\mid n}\ind_{\rm bin}
 R_z(n\pm hdu)\ind_{\rho_{j,z}(n,d)\le KW_1/L}\le KW_1=\mathfrak b.
\end{equation}
On a closed word the product of the endpoint rough indicators equals the product of the arrival indicators,
with visit multiplicities. This rewriting has already been made before closure is dropped.
Retain the arrival rough factor in every $c_i$ while dropping block and the other endpoint restrictions.
Set $F_{m+1}=1$ and recursively define
\begin{equation}\label{eq:rough-padding-recursion}
\begin{aligned}
 F_i(n)&=\sum_u c_i(n,u)F_{i+1}(n+\epsilon_i h d_i u)
 \le KW_1\|F_{i+1}\|_\infty,\\
 F_i(n)&\le\mathfrak b^{m-i+1},\qquad
 F_1(n)\le(KW_1)^m,\qquad m=2k.
\end{aligned}
\end{equation}
The eligible-degree condition may be retained or dropped in this positive row bound; either choice gives the
same upper bound. Backward induction therefore gives padding mass at most $(KW_1)^{2k}$.
This works for every $\mathfrak b\ge0$, including $\mathfrak b<1$, and survives singleton overwrites.

There are at most $Jk+S_1/2$ distinct centre labels. As in [O, (7.6)--(7.7), p.~38], their remaining reciprocal
sum and the forest count give
\begin{equation}\label{eq:moment-completion}
 \mathfrak b^{2k}e^{C_4Jk}(2W)^{Jk+L^{1/4}/2}.
\end{equation}
Signs, designations, $2^{S_1}$ and the external $e^{108L}$ are absorbed in $e^{C_4Jk}$.
For fixed $W$, $(2W)^{L^{1/4}/2}\le e^k$ eventually.
The logarithm of the resulting main bound is at least $-O_W(L\log L)$ because $W_1\gg L^{-1/6}$.
Both the exceptional classes and finite-comparison error $e^{-L^9}$ are smaller; one absolute enlargement of
$C_2$, before choosing $W$, absorbs them. This proves \eqref{eq:masked-moment}.

Finally redo [O, Lemma~8.2 and Proposition~8.3, pp.~40--41]. The absolute centred factor $a_d$ is constant along
a $d$-edge. For an outgoing increasing edge from $n$ to $n+hdu$, the weighted row reads
$R_z(n+hdu)$ and is bounded by $L\rho_{j,z}^+(n,d)$. For an incoming edge, read it from its larger endpoint
$n$ to $n-hdu$; $u\mid n$ is equivalent to divisibility at the lower endpoint, and the row is bounded by
$L\rho_{j,z}^-(n,d)$. Since $\rho_{j,z}=\max(\rho_{j,z}^+,\rho_{j,z}^-)$, each orientation contributes at most
$KW_1a_d(n)$. Weighted Schur with weight $\sqrt w$ and \eqref{eq:density-row} gives
\begin{equation}\label{eq:masked-schur}
 \|A_d\|\le2\mathfrak b,\quad
 \|A_dv\|^2\le4\mathfrak b^2\sum_n a_d(n)^2|v_n|^2,\quad
 \sum_d\|A_dv\|^2\le4\mathfrak b^2(8W)^J\|v\|^2\quad(v=Ev).
\end{equation}
Indeed $\sum_da_d(n)^2\le\prod_i(\omega_{\Pset_i}(n)+V_i)\le(8W)^J$ on $E$, by the arithmetic--geometric
mean inequality. No commutation with $E$ is used.
Since $\rho(H)^{2k}\le\|H^k\|_{\HS}^2$, Markov gives
$\rho(H)\le e\mathfrak b(C_2\sqrt W)^J$ outside $e^{-2k}$ of the origins.
Input~\ref{input:transfer} now gives \eqref{eq:masked-spectrum} with, for example,
$C\ge3\max(2,2\sqrt8,eC_2)$. All hypotheses have nonnegative row constants; none requires $\mathfrak b\ge1$.
\end{proof}

\section{Centering and the hard rough theorem}\label{sec:rough-proof}

\begin{proposition}[Polynomial on rough shifts]\label{lem:shift-polynomial}
Fix $a<199/200$, and suppose every element of $\mathcal Z\subset[D,2D)$ avoids the primes up to $e^{L^a}$,
where $H_0/2\le D\le e^{2L}$. For arbitrary $|c_b|=1$, put
\begin{equation}\label{eq:shift-polynomial}
 \mathcal B(\theta)=\sum_{b\in\mathcal Z}\frac{c_b\e(hb\theta)}b.
 \qquad \|\mathcal B\|_\infty\ll L^{-a},\qquad
 \int_0^1|\mathcal B(\theta)|^4d\theta\ll D^{-1}L^{-4a}.
\end{equation}
\end{proposition}
\begin{proof}
We repeat \cite[Lemma~4.2, pp.~13--14]{OpenAI007} with the changed cutoff $e^{L^a}$.
For the sup norm, count one rough variable in $[D,2D)$.
Orthogonality in the fourth moment imposes $b_1+b_2=b_3+b_4$; after discarding the unimodular coefficients and
using weights at most $D^{-4}$, count the four forms $x_1,x_2,x_3,x_1+x_2-x_3$ in a three-dimensional box of
side $D$. The last form is positive there.
Their bad-prime union densities are $1/p$ and $4/p+O(p^{-2})$, respectively. Every pair of the four defining
hyperplanes is independent, including at $2$, and the vector $(1,1,1)$ avoids all four; their union density is
strictly below one.

Use even inclusion--exclusion of order $r_0=2\lceil500\log L\rceil$.
For each selected prime set of product $q$, CRT in the interval or box gives normalized error $O(q/D)$.
There are at most $(r_0+1)(1+\pi(e^{L^a}))^{r_0}$ terms and $q\le e^{r_0L^a}$, so the total error is
\begin{equation}\label{eq:shift-crt-error}
 D^{-1}\exp(O(L^a\log L)).
\end{equation}
It is smaller than every fixed power of $L^{-1}$ because $a<199/200$.
The independent truncation remainder is at most
$(\sum_{p\le e^{L^a}}\nu(p))^{r_0+1}/(r_0+1)!\ll L^{-100}$, since that sum is at most
$(4+o(1))\log L$. The exact products are $O(L^{-a})$ and $O(L^{-4a})$.
Thus the two counts are $O(DL^{-a})$ and $O(D^3L^{-4a})$, giving \eqref{eq:shift-polynomial}.
\end{proof}

\begin{proposition}[Density in the testing and overlap step]\label{prop:testing}
Use the rough matrix \eqref{eq:rough-matrix}, an origin interval of length at least $e^{L^{625}/2}$, and the vector
$\sqrt w\lambda R_z$, truncated to its low-$\Qset$ state and then projected by $E$.
Its averaged squared norm is $O_h(MSW_1)$.
For a bin the retained centred sum has bound
\begin{equation}\label{eq:bin-testing}
 |C_j^\circ|\ll_h \frac{SW_1\,KW_1(C\sqrt W)^J}{L}+e^{-L^{0.9}}.
\end{equation}
After summing bins and dividing by $S_0$, the spectral contribution is
\begin{equation}\label{eq:rough-spectral-contribution}
 O_h\left(W_1^2\frac K\eta(C/\sqrt W)^J\right)+e^{-L^{0.8}}.
\end{equation}
\end{proposition}
\begin{proof}
Repeat \cite[Proposition~8.4, pp.~41--42]{OpenAI007} with all density factors retained.
The squared vector is the truncated weight $wR_z$. Its product-law mean is at most $SW_1$, since the small-prime
coordinates are disjoint from $\Qset$. Proposition~\ref{prop:threshold} transfers this mean to each of the $M$
block positions, giving the norm bound. Test \eqref{eq:masked-spectrum} on the nonexceptional blocks.
On every block the weighted absolute form is at most a fixed power of $L$ times $M$: use
$w\le17^{18\log L}$ and $\sum_da_d\le(8W)^J$ on the exterior projection. Thus the fraction $e^{-2k}$ of
exceptional blocks costs exponentially little even after division by $SW_1^2V_*$.

For completeness, the overlapping-block identity of [O, (8.12)--(8.13), p.~42] is unchanged.
If an increasing edge has displacement $q'=hdu$, the number of blocks among $t=0,\ldots,N-1$, $N=\lfloor T_j\rfloor$,
containing it is
\begin{equation}\label{eq:overlap}
 m_N(n,q')=[\min(N-1,n-1)-\max(0,n+q'-M)+1]_+.
\end{equation}
For $M\le n\le N$ this is $M-q'$. The boundary involves at most $2M$ sites, and the remaining discrepancy is
bounded by the deterministic row bound times
\begin{equation}\label{eq:overlap-error}
 O_h\left(\frac{he^{100L+1}}M+\frac M{T_j}+\frac1{T_j}\right).
\end{equation}
The row bound divided by $M$ is a fixed power of $L$, $M=e^{103L}+O(1)$, and $T_j\ge e^{L^{625}}$.
This includes floor errors and edges passing the prefix endpoint. Dividing the quadratic form by $2LM$ gives
\eqref{eq:bin-testing}. There are $O(L/\eta)$ bins and $S_0\ge SV_*/2$, $V_*\ge W^J$; these give
\eqref{eq:rough-spectral-contribution}. The argument did not use a signed-norm monotonicity principle.
\end{proof}

\begin{proof}[Proof of Theorem~\ref{thm:rough}]
Choose a permissible absolute $C\ge4$ in Proposition~\ref{prop:masked-moment} before $W$, and set
\begin{equation}\label{eq:rough-parameters}
\begin{gathered}
 A=625,\quad L=(\log X)^{1/A},\quad\beta=1/6,\\
 W=4eC^2,\quad\alpha=1/(1200W),\quad J=\lfloor\alpha\log L\rfloor,\\
 \eps=\alpha/100,\quad\eta=L^{-\eps},\quad K=L^{\alpha\log2+\eps}.
\end{gathered}
\end{equation}
All labels exceed $z$. The function $f_z=\lambda R_z$ is completely multiplicative and, for a graph label $q'=du$,
\begin{equation}\label{eq:rough-dilation}
 f_z(q'm)f_z(q'(m+h))=f_z(m)f_z(m+h).
\end{equation}
The square $f_z(q')^2=1$ proves this without a coprimality assumption between $q'$ and $m$.
The full-divisibility part of the centred sum is exactly
\begin{equation}\label{eq:raw-rough}
 \sum_j\sum_{(d,u)\in\mathcal S_j}\frac{16^{\omega(u)}}{du}
 \frac1{T_j/(du)}\sum_{m\le T_j/(du)}f_z(m)f_z(m+h).
\end{equation}
Its cutoff is in $(Xe^{-\eta},X]$. The positive upper sieve bounds the missing interval by
$O_h(\eta XV_h(z))$ and the full rough count by $O_h(XV_h(z))$.
If an actual tail is shorter, enlarge it to the complete interval $(Xe^{-\eta},X]$.
The required sieve moduli are $X^{o(1)}$, so this bound is uniform in $z$.
Thus \eqref{eq:raw-rough} differs from $S_0X^{-1}\sum_{n\le X}f_z(n)f_z(n+h)$ by
$O_h(S_0\eta V_h(z))+e^{-L}$.

The padding deletion is Proposition~\ref{prop:collision}.
For the degree deletions repeat \cite[Lemma~4.6, pp.~19--20]{OpenAI007}, verifying the changed tilted law.
Under $w/S$, the probabilities are $17/(p+16)$, and exponential Markov with parameter $1/20$ gives
\begin{equation}\label{eq:q-degree-tail}
 \Prob^{\rm tilt}(\omega_{\Qset}>18\log L)
 \ll_h L^{-18/20+17(e^{1/20}-1)}\ll_h L^{-1/40}.
\end{equation}
The centre tilt by $a_d^+$ forces at most $J$ primes and leaves the others independent, of total mean at most
$2WJ$. Its tail is at most
$\exp(-6WJ+J+2(e-1)WJ)\le e^{-2WJ}$.
In each fixed numerical label the small-prime pair has density $V_h(z)$; only these coordinates are independent
of the centre and padding degree events. This retains the pair density in both deletions.
For rare sites, the second moment of
$w\prod_i(\omega_{\Pset_i}+V_i)$ is a fixed power of $L$: its $\Qset$ factor is
$\prod_{p\in\Qset}(1+288/p)$, and the centre factors have moments bounded by fixed multiples of $V_i^2+V_i$.
Cauchy--Schwarz with \eqref{eq:rare}, followed by the $O(L/\eta)$ bins, gives exponentially small error.
All these are positive finite-residue expansions covered by Proposition~\ref{prop:threshold}; the padding
threshold is encoded only after the low-degree cutoff.
Together with Proposition~\ref{prop:testing}, the non-centering errors, divided by $V_h(z)$, are
\begin{equation}\label{eq:rough-budget}
 C_h\left\{\eta+\frac{2^J(1+\log L)^6}{K}
 +2^J(L^{-1/40}+e^{-2WJ})+\frac K\eta(C/\sqrt W)^J\right\}+e^{-L^{0.8}}.
\end{equation}

It remains to redo \cite[Lemmas~4.1--4.2, 4.4, pp.~13--16]{OpenAI007} with $f_z$, the new range, and the densities.
Use Proposition~\ref{lem:shift-polynomial} with
$a=9949/10000$, and set $\Gamma=89/100$, $s=17/100$.
In a partial centred term, a nonempty subset $I$ of the $J$ supplies is replaced by $-1/p$.
Writing $b=\prod_{i\in I}p_i$, $d'=d/b$, and $n=ud'v$ extracts a reciprocal mass at most $SV_*$ and leaves
\begin{equation}\label{eq:partial-rough}
 \frac1Y\sum_{v\le Y}\sum_{b\in\mathcal Z}\frac{c_b}b f_z(v)f_z(v+hb),\qquad Y=T_j/(ud').
\end{equation}
The allowed $b$ lie in at most two dyadic intervals, $H_0/2\le D\le e^{2L}$, and $Y\ge e^{L^{625}/2}$.
The coefficients are unimodular. This is [O, (4.8), p.~16] with the exact identity
\eqref{eq:rough-dilation} in place of its original Liouville identity.

Replacing $\lambda(p)$ by zero at $p\le z$ reduces \eqref{eq:pretentious} by at most
$\sum_{p\le z}1/p\le\beta\log L+O(1)$.
Input~\ref{input:mrt} therefore gives $M(f_z;Y,Q)\ge(625/5)\log L$ eventually.
Here $10\le D\le Y$ and $Q$ has exactly the prescribed cap; the three errors in \eqref{eq:mrt} are bounded by
\begin{equation}\label{eq:rough-mrt-errors}
 L^{-25/4},\qquad L^{-199/200}\log L,\qquad L^{-25/28},
 \qquad \frac1Y\sum_{v\le Y}\left|\sum_{m=1}^D f_z(v+m)e(\theta m)\right|\ll D L^{-\Gamma}.
\end{equation}
The last assertion follows from unit intervals of origins, as in [O, (4.6), p.~15]; endpoint cost $O(D/Y)$ is
exponentially small.

Average \eqref{eq:partial-rough} over translations $1\le m\le D$; the endpoint cost is $O_h(D/Y)$.
Fourier orthogonality gives $(YD)^{-1}\sum_{v\le Y}\int\mathcal B F_vG_v$, with
$F_v=\sum_{m=1}^D f_z(v+m)e(\theta m)$ and
$G_v=\sum_{a=1}^{(2h+1)D}f_z(v+a)e(-\theta a)$.
For any such window, the elementary upper sieve gives rough count $O_h(DW_1)$.
For clarity, even inclusion--exclusion of order $O(\log L)$ has total CRT endpoint error
$D^{-1}\exp(O(L^{1/6}\log L))$, and its independent remainder can be any fixed negative power of $L$.
This verifies the uniformity even for the smallest $D$.
Thus $\|F_v\|_2\|G_v\|_2\ll_h DW_1$ and $|G_v|\ll_h DW_1$.
On $|\mathcal B|\le L^{-1-s}$ use Parseval; on the complement its measure is at most
$L^{4+4s}\int|\mathcal B|^4\ll D^{-1}L^{4+4s-4a}$.
Use also $\|\mathcal B\|_\infty\ll L^{-a}$ and \eqref{eq:rough-mrt-errors}. The two partial contributions are
\begin{equation}\label{eq:rough-centering-pieces}
 O_h(W_1L^{-1-s}),\qquad O_h(W_1L^{4+4s-5a-\Gamma}).
\end{equation}
After at most $2^J$ subsets, $O(L/\eta)$ bins and division by $V_h\asymp_h W_1^2$, their exponents are
\begin{equation}\label{eq:rough-centering-exponents}
 \beta+\eps+\alpha\log2-s,\qquad
 \beta+\eps+\alpha\log2+5(1-a)+4s-\Gamma.
\end{equation}
Without the $\eps,\alpha$ terms the margins are $-1/300$ and $-107/6000$.
For $W\ge10$ both are less than $-\eps$.

In \eqref{eq:rough-budget}, the main deletion is $L^{-\eps}(1+\log L)^6$ and, since
$C/\sqrt W=1/(2\sqrt e)$, the spectral exponent is $-\alpha/2+2\eps$.
The remaining degree terms have strict negative margins. Absorbing logarithmic factors gives
\begin{equation}\label{eq:rough-joint}
 \left|\sum_{n\le X}\lambda(n)\lambda(n+h)R_z(n)R_z(n+h)\right|
 \ll_h XV_h(z)L^{-\eps/2}.
\end{equation}
Lemma~\ref{lem:rough-sieve} supplies the unsigned and both single-sign terms with arbitrary-logarithmic absolute
errors. Choose their saving larger than $2/3750+\gamma$ and use the exact sign expansion
\begin{equation}\label{eq:sign-expansion}
 \ind_{\lambda(n)=s,\lambda(n+h)=t}
 =\tfrac14(1+s\lambda(n))(1+t\lambda(n+h)).
\end{equation}
Equation~\eqref{eq:rough-main} follows, for example with $\gamma=\eps/(4\cdot625)$.
The cutoffs satisfy $\log z\le L^{1/6}$, exactly the stated range. All graph constants precede $W$ and $X$;
choose the final threshold after $h$ and those constants are fixed. Origin exceptions were internal averages,
not exclusions of values of $X$.
\end{proof}

\section{Balanced soft small-prime weights}\label{sec:soft}

\begin{lemma}[Prime-square masks and densities]\label{lem:soft-density}
For $0\le r_1,r_2\le1$ the function
\begin{equation}\label{eq:two-soft-mask}
 A_z(n)=S_z(n)S_z(n+h)r_1^{\omega_z(n)}r_2^{\omega_z(n+h)}
\end{equation}
is an average of negative prime-divisibility CNF indicators, with total averaging mass one, together with the
always retained negative prime-square tests. Its unsigned mean is compared by Lemma~\ref{lem:cnf} in the prime-square
model. When $r_1=r_2=r$, its pair mean is exactly \eqref{eq:soft-local}. Its single and squared means are
\begin{equation}\label{eq:soft-single-products}
 W_r=\prod_{p\le z}(1-1/p+r(p-1)/p^2),\qquad
 W_{r^2}=\prod_{p\le z}(1-1/p+r^2(p-1)/p^2).
\end{equation}
Uniformly for $0\le r\le1$, $\delta=1-r$, and $z\ge2$,
\begin{align}
 \log W_r&=-\delta\log\log(2z)+O(1),\label{eq:soft-log-first}\\
 \log W_{r^2}&=-(2\delta-\delta^2)\log\log(2z)+O(1),\label{eq:soft-log-square}\\
 \log V^{\sfw}_{r,h}(z)&=-2\delta\log\log(2z)+O_h(1).\label{eq:soft-log-pair}
\end{align}
In particular
\begin{equation}\label{eq:soft-density-ratios}
\begin{gathered}
 W_r\asymp v^{-\delta},\quad W_{r^2}\asymp v^{-(2\delta-\delta^2)},\quad
 V^{\sfw}_{r,h}\asymp_h v^{-2\delta}\asymp_h W_r^2,\\
 W_{r^2}/V^{\sfw}_{r,h}\ll_h v^{\delta^2},\quad W_r/V^{\sfw}_{r,h}\ll_h v^\delta,
 \qquad v=\log(2z).
\end{gathered}
\end{equation}
\end{lemma}
\begin{proof}
Independently retain $p\nmid n$ with probability $1-r_1$, and $p\nmid n+h$ with probability $1-r_2$.
Always retain $p^2\nmid n,n+h$. On a coordinate modulo $p^2$, $p\nmid n$ excludes the $p$ residues
$0,p,\ldots,(p-1)p$, and a negative square test excludes one residue. Each realization is thus a conjunction
of at most $2\sum_{p\le z}(p+1)$ negative equalities. For $z\le e^{L^b}$, $b<1/2$, the clause count is less than
$e^{L^3}$ and the moduli are at most $e^L$. Averaging Lemma~\ref{lem:cnf} costs no scalar coefficient mass.
The same argument is valid with $b<1$ whenever the coordinate-size condition holds separately.

List the residues modulo $p^2$. Outside the forbidden roots modulo $p$ the pair weight is one.
A root at one endpoint has $p-1$ admissible lifts of weight $r$.
A common root has $p-1$ or $p-2$ admissible lifts of weight $r^2$, according as $p^2\mid h$ or not.
This gives \eqref{eq:soft-local}; at $p=2$, $h\equiv2\pmod4$, the factor is exactly $1/2$.
The analogous single-site count gives \eqref{eq:soft-single-products}.

For the full range $0\le r\le1$ the single local factors are
\begin{equation}\label{eq:soft-local-expanded}
 q_r(p)=1-\delta/p-r/p^2,\quad
 q_{r^2}(p)=1-(2\delta-\delta^2)/p-r^2/p^2,\quad
 v_{r,h}(p)=1-2\delta/p-2r/p^2\quad(p\nmid h).
\end{equation}
The single factors are at least $1-1/p$. For $p\ge5$, the factors are at least $4/5,4/5,3/5$, respectively.
Taylor expansion of their logarithms therefore has a uniform $O(p^{-2})$ error.
The factors at $p=2,3$ and at $p\mid h$ are finitely many positive exceptions with uniform positive lower bounds:
the common-root factors are at least $1-1/p$, and the generic factor at $3$ is at least $1/3$.
Mertens' harmonic sum and $\sum_pp^{-2}<\infty$ give
\eqref{eq:soft-log-first}--\eqref{eq:soft-log-pair}; exponentiation gives \eqref{eq:soft-density-ratios}.
Only these product comparisons are uniform in $r$. Theorem~\ref{thm:soft} still fixes $r$ before $X$.
\end{proof}

\begin{proposition}[Classical single-sign orthogonality]\label{prop:soft-single}
For each fixed $0<\rho_*<1$ and every $B>0$, uniformly for
$2\le z\le\exp((\log X)^{\rho_*})$ and $r_1,r_2\in[0,1]$,
\begin{equation}\label{eq:soft-single-cancel}
 \sum_{n\le X}\lambda(n)A_z(n),\quad\sum_{n\le X}\lambda(n+h)A_z(n)
 \ll_{h,\rho_*,B}X(\log X)^{-B}.
\end{equation}
\end{proposition}
\begin{proof}
Average the negative prime tests in Lemma~\ref{lem:soft-density}. It suffices to treat uniformly the deterministic mask
$M(n)=\prod_{p\in\mathcal T_1}\ind_{p\nmid n}\prod_{p\in\mathcal T_2}\ind_{p\nmid n+h}$,
$\mathcal T_i\subset\{p\le z\}$, with the two squarefree masks still present.
Put $T=X^{1/100}$ and
\begin{equation}\label{eq:square-truncation}
 S_{z,T}(n)=\sum_{\substack{b\le T,\ P^+(b)\le z\\b^2\mid n}}\mu(b).
\end{equation}
The full identity is the same sum without $b\le T$.
Thus $|S_z-S_{z,T}|\le\sum_{b>T,b^2\mid n}1$ and $|S_{z,T}(n)|\le\tau(n)$.
Replacing both factors has $L^1$ error at most
\begin{equation}\label{eq:square-truncation-error}
 \sum_{n\le X}\sum_{b>T,b^2\mid n}1+
 \sum_{n\le X}\tau(n)\sum_{c>T,c^2\mid n+h}1
 \ll_\xi X/T+X^{1+\xi}/T=O_h(X^{199/200}),\quad \xi=1/200.
\end{equation}
Use the complete shifted prefix through $X+h$ and $\tau(n)\ll_\xi n^\xi$; no additional floor sum is charged.

For fixed $b,c\le T$, the two square divisibilities are inconsistent or give one class modulo
$q_0=[b^2,c^2]\le T^4$.
Tests at primes dividing $q_0$ are automatic or make the mask zero. Remove them.
At each remaining prime at most two roots are excluded, and CRT works with $d$ coprime to $q_0$.
Input~\ref{input:fl}, at level $D=X^{1/8}$, gives the $L^1$ upper-sieve approximation error on this progression
\begin{equation}\label{eq:soft-sieve-error}
 O\left(\frac X{q_0}E(s)+D\log(2D)\right),\qquad
 s\ge\tfrac18(\log X)^{1-\rho_*}.
\end{equation}
Its local product is at most one; its dimension condition is uniformly bounded by the two-root sieve.
The modulus and prime masks introduce no new constants. Summing this error costs
$O(XE(s)+T^2D\log(2D))$, since
\begin{equation}\label{eq:square-lcm}
 \sum_{b,c\ge1}[b^2,c^2]^{-1}\le\zeta(2)^3,
 \qquad T^2D=X^{29/200}.
\end{equation}
For the inequality write $b=gu,c=gv$, $(u,v)=1$ and drop that coprimality.
Inserting a Liouville sign cannot increase this $L^1$ error.

All remaining sums have modulus $q=q_0d\le X^{33/200}$.
For fixed $q$ there are at most $\tau(q)^3$ choices of $b,c,d$ and at most $\tau(q)$ root classes per choice.
The absolute remainder is bounded by
\begin{equation}\label{eq:soft-grouped-bv}
 \sum_{q\le X^{33/200}}\tau(q)^4E_q,
 \quad E_q=\max_{a\bmod q}\max_{y\le2X}\left|\sum_{\substack{n\le y\\n\equiv a\pmod q}}\lambda(n)\right|.
\end{equation}
Lemma~\ref{lem:bv} gives $\sum E_q\ll_A X\log^{-A}X$ and $E_q\le2X/q+1$.
Since $\tau(q)^8\le\tau_{256}(q)$, the divisor-variable bounds give
\begin{equation}\label{eq:soft-bv-cauchy}
 \sum_{q\le X^{33/200}}\tau(q)^8E_q\ll X\log^{256}(2X),\qquad
 \sum\tau(q)^4E_q\ll_A X\log^{128-A/2}X.
\end{equation}
Choose $A$ sufficiently large. The translated sign is a difference of two prefixes in another residue class.
This proves \eqref{eq:soft-single-cancel} using only the classical FL and BV inputs.
\end{proof}


\begin{lemma}[Soft padding recursion]\label{lem:soft-recursion}
Fix $0<r<1$, use padding $16$, and let $b=b_z$ and $b^+$ be its cap in \eqref{eq:soft-cap}.
In a fixed non-$\Pset$ residue environment put
\begin{equation}\label{eq:soft-recursion-probabilities}
 \rho_j^\pm(n,d)=\frac1{w(n)}\sum_{\substack{u\mid n,\ u\in\Qset^{\sfw}\\j\eta\le\log(du)<(j+1)\eta}}
 a(u)b^+(n\pm hdu),\qquad \rho_j=\max(\rho_j^+,\rho_j^-).
\end{equation}
For fixed $d_i$ and $\epsilon_i\in\{\pm1\}$ define
\begin{equation}\label{eq:soft-recursion-coefficients}
 c_i^0(n,u)=\frac{La(u)}{w(n)}\ind_{u\mid n}\ind_{\rm bin}
 \ind_{\rho_j(n,d_i)\le KW_r/L}.
\end{equation}
With $F_{m+1}=1$, the recursion obeys
\begin{equation}\label{eq:soft-padding-recursion}
\begin{aligned}
 F_i(n)&=\sum_u c_i^0(n,u)b(n+\epsilon_i h d_i u)F_{i+1}(n+\epsilon_i h d_i u)\\
 &\le KW_r\|F_{i+1}\|_\infty,\qquad F_1(n)\le(KW_r)^m.
\end{aligned}
\end{equation}
On a closed word the balanced edge factors satisfy
\begin{equation}\label{eq:soft-closed-weight}
 \prod_{i=1}^m\sqrt{b(n_i)b(n_{i+1})}
 =\prod_{i=1}^m b(n_{i+1})=\prod_{i=1}^m b(n_i),\qquad n_{m+1}=n_1,
\end{equation}
including repeated visits and zero weights.
\end{lemma}
\begin{proof}
Since $0\le b\le b^+$, the row sum of $c_i^0(n,u)b(n+\epsilon_i h d_i u)$ is at most
$L\rho_j^{\epsilon_i}(n,d_i)\ind_{\rho_j\le KW_r/L}\le KW_r$.
Here $\rho_j^{\epsilon_i}$ means $\rho_j^+$ for $\epsilon_i=1$ and $\rho_j^-$ for $\epsilon_i=-1$.
This is the backward induction of \cite[Lemma~7.3, p.~37]{OpenAI007} with a nonnegative row constant $KW_r$.
It is uniform in the starting site and does not require $KW_r\ge1$.
For \eqref{eq:soft-closed-weight}, every visited factor occurs in the two square roots at its incident edges,
counted with visit multiplicity. Nonnegativity gives the identity, including when a factor is zero.
Make it before dropping closure in the positive path enlargement; then \eqref{eq:soft-padding-recursion} at $m=2k$
gives $(KW_r)^{2k}$. Adding the eligible padding-degree condition only decreases these sums.
\end{proof}

\begin{proposition}[Balanced operator and capped cutoff]\label{prop:soft-graph}
Fix $0<r<1$ and retain \eqref{eq:rough-parameters}. Let $b=b_z$.
Use the symmetric edge matrix \eqref{eq:soft-matrix} below and test it with $\sqrt{w(n)b(n)}\lambda(n)$.
Define
\begin{equation}\label{eq:soft-cap}
 M_r=\left\lceil\frac{100\log L}{-\log r}\right\rceil,\qquad
 b^+(n)=\max(b(n),r^{M_r}),\qquad r^{M_r}\le L^{-100}.
\end{equation}
Use $b^+(n\pm hdu)$ in the target padding probabilities \eqref{eq:masked-rho}, take their maximum $\rho_j$, and
cut at $\rho_j\le KW_r/L$, after the low-$\Qset$ cutoff. The actual matrices and vectors use $b$. More explicitly, for $m=n+hdu$ let
\begin{equation}\label{eq:soft-matrix}
\begin{aligned}
 \tau_j^b(n,d)&=\ind_{\omega_{\Qset}(n)\le\lfloor18\log L\rfloor}
 \ind_{\rho_j(n,d)\le KW_r/L}\sigma_j(n),\\
 A_d^b(n,m)&=\frac{L16^{\omega(u)}}{\sqrt{w(n)w(m)}}\ind_{u\mid n}c_d(n)
 \tau_j^b(n,d)\tau_j^b(m,d)\sqrt{b(n)b(m)}.
\end{aligned}
\end{equation}
Set $A_d^b(m,n)=A_d^b(n,m)$ and other entries to zero. In the following proof write $A_d=A_d^b$.
The associated centred sums use $f=g=\lambda b$, $a=16$, $\epsilon=1$ and $\tau_j=\tau_j^b$.
They satisfy the finite comparison of Proposition~\ref{prop:threshold}, the moment and spectral bounds of
Proposition~\ref{prop:masked-moment} with $W_1$ replaced by $W_r$, and averaged testing norm $O_{h,r}(MSW_r)$.
Explicitly, with $H_b{}_{d,d'}=\ind_{d\ne d'}A_{d'}^b$,
\begin{equation}\label{eq:soft-moment-spectrum}
 \E_*\|H_b^k\|_{\HS}^2\le\{KW_r(C_2\sqrt W)^J\}^{2k},\qquad
 \left\|E\Bigl(\sum_d A_d^b\Bigr)E\right\|\le KW_r(C\sqrt W)^J.
\end{equation}
The spectral estimate excludes at most a fraction $e^{-2k}$ of the block origins; the finite-origin and
constant conventions are those of Proposition~\ref{prop:masked-moment}.
Their padding deletion is
\begin{equation}\label{eq:soft-deletion}
 O_{h,r}\left(W_r^2\frac{2^J(1+\log L)^6}{K}\right).
\end{equation}
\end{proposition}
\begin{proof}
The exact quadratic pairing contains $b(n)b(m)$, because each endpoint contributes a square root in the matrix
and a square root in the vector. At zero-weight vertices all entries and the vector vanish.
On the remaining space weighted Schur uses the positive weight $\sqrt{wb}$ and gives
\begin{equation}\label{eq:soft-balancing}
 |A_d(n,m)|\frac{\sqrt{w(m)b(m)}}{\sqrt{w(n)b(n)}}
 =\frac{L16^{\omega(u)}}{w(n)}a_d(n)\ind_{u\mid n}b(m)\times\text{retained cuts}.
\end{equation}
Thus the oriented padding row is at most $KW_r$.
On a closed word,
\begin{equation}\label{eq:soft-arrival}
 \prod_{i=1}^{2k}\sqrt{b(n_i)b(n_{i+1})}=\prod_{i=1}^{2k}b(n_{i+1}),\qquad n_{2k+1}=n_1.
\end{equation}
This includes repeated visits and zeros. The denominators have the identical closed-word identity for $w$.
Apply it before dropping closure; no square root remains to be encoded.

For finite comparison, fix an active state of at most $18\log L$ padding primes, hence at most
$N\le L^{18\log2}$ candidates. A bad weighted sum is the OR over level vectors
$(k_1,\ldots,k_N)\in\{0,\ldots,M_r\}^N$ whose scalar sum exceeds the threshold, of the events
$b^+(n_i)\ge r^{k_i}$ for all $i$. If $k_i=M_r$ this event is automatic; otherwise it is exactly
\begin{equation}\label{eq:cap-level}
 S_z(n_i)=1\quad\text{and}\quad\omega_z(n_i)\le k_i.
\end{equation}
Negation yields an AND of ORs of square divisibilities $p^2\mid n_i$ or conjunctions of $k_i+1$ distinct prime
divisibilities. The square-divisibility alternative is essential at $b=0$.
On a $p^2$ coordinate, $p\mid n_i$ is an OR of $p$ equalities and $p^2\mid n_i$ one equality.
Substitute the interval DNFs of Proposition~\ref{prop:threshold}, distribute the bottom $O_r(\log L)$
conjunctions and flatten ORs. The depth is three, and the logarithmic size is at most
\begin{equation}\label{eq:soft-circuit-size}
 N\log(M_r+1)+O_r(L^{1/6}\log L+\log^2L)+O(L\log L)<L^{14}.
\end{equation}
The coordinate moduli are at most $e^L$, pairwise coprime to the disjoint graph coordinates.
Use $t=\lceil L^{589}\rceil$ and the identical CRT/Walsh correction, giving aggregate error $e^{-L^9}$.
Every residual factor $b$ is the average in Lemma~\ref{lem:soft-density}; repeated occurrences use independent
auxiliary marks, which correctly produce $b$ to that multiplicity. Their probabilities sum to one.
They only add CNF conditions and introduce no scalar-mass cost or physical-site independence.

The padding experiment is unchanged from Proposition~\ref{prop:collision}.
For different eligible paddings the generic three-site factor, at $p>3$, $p\nmid h(u_1-u_2)$, is exactly
$1-3\delta/p-3r/p^2$. Comparing it with three single-site factors changes a convergent $O(p^{-2})$ product.
Lemma~\ref{lem:density} bounds exceptional primes as before, and gives
\begin{equation}\label{eq:soft-triple}
 \E_{\rm small}[b(n)b(n\pm hdu_1)b(n\pm hdu_2)]
 \ll_{h,r}W_r^3(1+\log L)^6.
\end{equation}
At equal paddings the mean of $b(n)b(n+hdu)^2$ has generic logarithmic coefficient
$-(3\delta-\delta^2)/p$ and is at most $C_{h,r}W_r^3v^{\delta^2}$.
If $\rho_j^{(b)}$ denotes the uncapped target probability, then
\begin{equation}\label{eq:soft-second}
 \E_*[wb\sum_j(\rho_j^{(b)})^2]
 \ll_{h,r}SW_r^3\{(1+\log L)^6/L+v^{\delta^2}(v/L)^{24/17}\}
 \ll_{h,r}SW_r^3(1+\log L)^6/L.
\end{equation}
The diagonal-to-main exponent is at most
$(1/6)(24/17+1)+1-24/17=-1/102$.
As before the different-padding estimate is used on eligible bin pairs before enlargement.
Since $b^+\le b+r^{M_r}$, the cap's added squared collision is
$O_r(L^{-200}SW_r/L)$, by the unmasked collision estimate and independence of the small-prime coordinates.
Lemma~\ref{lem:soft-density} gives $W_r\gg L^{-1/6}$ for this whole fixed-$r$ range, so the extra term is negligible.
The maximum of the two orientations only costs a constant.
On a bad row $wb\rho_j^{(b)}\le Lwb\rho_j^2/(KW_r)$. Multiply the independent centre majorant $2^J/d$,
sum, and obtain \eqref{eq:soft-deletion}.

We now check the changed graph proof in the same order as Proposition~\ref{prop:masked-moment}.
The trace expansion [O, pp.~24--25] uses \eqref{eq:soft-arrival} and remains a signed finite expansion.
In [O, Lemma~5.5, p.~26] the integer powers of $b$, the capped cutoff and the padding state belong to the
non-$\Pset$ factor $R$; singleton overwrites leave them unchanged, and $G$ is still only the witness factor.
The three negligible classes drop $b$ using $0\le b\le1$, only in their nonnegative majorants:
[O, Lemma~5.6, p.~26] retains the $U/2$ extra unlit reciprocals;
[O, Lemma~6.2, pp.~28--29] retains the high-rank lit-consistency equations;
[O, Lemma~6.5, pp.~31--33] retains the witness and padding divisibilities and the ordered unit-coefficient
eliminations. In each class the pointwise comparison \eqref{eq:crude-domination} gives the same
$\exp(O(L\log L))$ change of padding coefficient, and no $b$ is replaced by a positive lower bound.
After these three classes are discarded, the forest hypotheses [O, Lemmas~7.1--7.2, pp.~34--37] hold on the
smaller support. On surviving words keep the arrival $b$: Lemma~\ref{lem:soft-recursion}, with the closed identity
\eqref{eq:soft-closed-weight} made before closure is dropped, gives the padding bound $(KW_r)^{2k}$.
The completion [O, pp.~37--38] retains the same $C_2$, since $W_r\gg L^{-1/6}$ makes the main logarithm at least
$-O_W(L\log L)$, even if $KW_r<1$.
Weighted Schur [O, Lemma~8.2, p.~40] uses \eqref{eq:soft-balancing} on $b>0$, on each invariant level set of
$a_d$. The outgoing row uses $\rho_j^+$ and the incoming row $\rho_j^-$, each capped at $KW_r/L$.
Vertices with $b=0$ have zero rows and columns, so no division by zero occurs. Input~\ref{input:transfer} applies
with $a=2KW_r$, $b=2KW_r(8W)^{J/2}$. This gives the same permissible absolute $C$.
Finally the truncated squared vector is $wb$ with independent mean at most $SW_r$.
The overlap identity \eqref{eq:overlap} and deterministic exceptional-block argument apply with the same fixed
polynomial bounds. This proves all asserted testing and spectral estimates.
\end{proof}

\begin{proof}[Proof of Theorem~\ref{thm:soft}]
By Lemma~\ref{lem:soft-density}, $V^{\sfw}_{r,h}\asymp_{h,r}W_r^2$ throughout the fixed-$r$ range.
For every label $q'=du$, whose primes exceed $z$,
\begin{equation}\label{eq:soft-dilation}
 b(q'm)=b(m),\quad (\lambda b)(q'm)=\lambda(q')(\lambda b)(m),\quad
 (\lambda b)(q'm)(\lambda b)(q'(m+h))=(\lambda b)(m)(\lambda b)(m+h).
\end{equation}
The small prime-square valuations are preserved even when $q'$ and $m$ share large primes.
Thus the raw identity and normalization are unchanged, with a pair density $V^{\sfw}_{r,h}$.
Positive window bounds for the pair and single masks follow by averaging the elementary upper sieve over the
marks of Lemma~\ref{lem:soft-density}, dropping square restrictions for an upper bound.
The local products change by bounded positive factors, since the difference is $O(p^{-2})$ outside fixed
exceptional primes. An $O(\log L)$ truncation has normalized CRT error
$H^{-1}e^{O(L^{1/6}\log L)}$ for $\log H\gg L^{199/200}$, and negligible independent remainder.
This proves bounds $O_{h,r}(HV^{\sfw}_{r,h})$, $O_r(HW_r)$ and $O_r(HW_{r^2})$, including the raw tail of
length $\eta X$.

Proposition~\ref{prop:soft-graph} gives the relative spectral term
\begin{equation}\label{eq:soft-spectral}
 O_{h,r}((K/\eta)(C/\sqrt W)^J).
\end{equation}
Its padding deletion is \eqref{eq:soft-deletion}; degree and rare-site deletions preserve the pair density by the
disjoint small prime-square coordinates. Their bounds are those in \eqref{eq:rough-budget}.
For centering $f=\lambda b$ is multiplicative and $1$-bounded; its prime values change from those of $\lambda$
only at $p\le z$ and by at most one. Thus \eqref{eq:rough-mrt-errors} still holds with $\Gamma=.89$.
The small-frequency piece uses the squared density $W_{r^2}$, and the large-frequency piece the pointwise mean
$W_r$. At both uses Lemma~\ref{lem:soft-density} supplies the full-$r$ comparisons.
The two relative exponents, with $a=.9949$, $s=.17$, $\beta=1/6$, are
\begin{equation}\label{eq:soft-centering-exponents}
 \beta\delta^2+\eps+\alpha\log2-s,\qquad
 \beta\delta+\eps+\alpha\log2+5(1-a)+4s-\Gamma.
\end{equation}
They do not exceed \eqref{eq:rough-centering-exponents}, because $\delta,\delta^2\le1$.
The raw tail has cost $O_{h,r}(\eta V^{\sfw}_{r,h})$.
The full relative budget therefore gives
\begin{equation}\label{eq:soft-joint}
 \left|\sum_{n\le X}\lambda(n)\lambda(n+h)b(n)b(n+h)\right|
 \ll_{h,r}XV^{\sfw}_{r,h}L^{-\eps/2}.
\end{equation}
The unsigned mass is its independent product plus an exponentially small finite-comparison error by
Lemma~\ref{lem:soft-density}; Proposition~\ref{prop:soft-single} gives arbitrary logarithmic savings for both
single signs. Expanding \eqref{eq:sign-expansion} proves \eqref{eq:soft-main}, with the same
$\gamma=\eps/(4\cdot625)$. The constants in the graph moments are absolute; circuit thresholds and positive
comparisons may depend on the fixed $r$. No choice $r=r(X)$ is made.
\end{proof}

\section{Bounds on the number of prime factors}\label{sec:consumer}

\begin{proposition}[Positive moments]\label{prop:positive-moment}
For fixed even $h$, $z\le X^{o(1)}$ and $1\le u\le2$,
\begin{align}
 \sum_{2\le n\le X}R_z(n)R_z(n+h)u^{\omega(n)}
 &\ll_h X(\log X)^{u-1}(\log z)^{-(u+1)}+O(\sqrt X\log^2X),\label{eq:hard-positive-moment}\\
 \sum_{n\le X}b_z(n)b_z(n+h)u^{\omega(n)}
 &\ll_h X(\log X)^{u-1}(\log z)^{-(u+1)\delta}+O(\sqrt X\log^2X),\label{eq:soft-positive-moment}\\
 \qquad \delta&=1-r.\notag
\end{align}
The same estimates hold at the other endpoint. The soft upper bound applies to every fixed $0<r<1$.
\end{proposition}
\begin{proof}
Apply the Nair--Tenenbaum theorem \cite[Theorem~1, pp.~1--2]{Henriot} on each dyadic interval $[Y,2Y]$ with
$Y\ge\sqrt X$, to $Q_1(t)=t$, $Q_2(t)=t+h$ and the nonnegative weights
\begin{equation}\label{eq:nt-functions}
 F(a,b)=u^{\omega(a)}R_z(a)R_z(b),\qquad
 F(a,b)=u^{\omega(a)}r^{\omega_z(a)+\omega_z(b)}.
\end{equation}
In the second bound the squarefree masks have been dropped.
The polynomials are primitive, irreducible and coprime; the product has degree two, discriminant $h^2\ne0$,
and no fixed prime divisor since $h$ is even.
For the function class in that theorem, coprime-product growth is bounded by $2^{\omega(a)}\le\tau(a)$, hence
both a fixed $A^{\Omega}$ bound and, for every fixed $\xi>0$, a uniform $B_\xi n^\xi$ bound.
Choose its interval exponent and coefficient exponent to be $1/2$, and then a permissible $\xi$.
The intervals have length $Y>Y^{1/2}$, and the coefficient threshold holds for large $Y$ with fixed $h$.
These verify the hypotheses uniformly in $z,u$ (and in the bounded prime values of $r$).

Both forms have one root modulo every integer. In the positive divisor sum on the right side of the cited
theorem, drop the product cutoff and enlarge to Euler factors over primes at most $2Y$.
For the hard weight, its local logarithms, including the theorem's sieve product, are
$(u-1)/p+O(p^{-2})$ for $p>z$, and $-2/p+O_h(p^{-2})$ for $p\le z$.
The fixed divisors of $h$ give bounded factors. This yields
$C_hY(\log Y)^{u-1}(\log z)^{-(u+1)}$.
For the soft weight the prime values are $ur,r$ below $z$ and $u,1$ above $z$, giving logarithm
\begin{equation}\label{eq:soft-moment-log}
 (u-1)\log\log X-(u+1)\delta\log\log z+O_h(1).
\end{equation}
Higher prime-power errors are uniformly summable. Geometric dyadic summation proves the stated main terms.
Below $\sqrt X$, $\sum2^{\omega(n)}\le\sum\tau(n)\ll\sqrt X\log X$, within the remainder.
Swapping the two functions proves the other endpoint. No signed or lower estimate is taken from this theorem.
\end{proof}

\begin{lemma}[A useful proportional tail]\label{lem:chernoff}
If a positive pair weight has mass of order $X\log^{-2\kappa}X$ and one-endpoint moments bounded by
$CX\log^{-2\kappa+(1-\kappa)(u-1)}X+O(\sqrt X\log^2X)$, $1\le u\le2$, then, for $0<\kappa<1/2$,
the mass with $\omega(n)>(1-\kappa/2)\log\log X$ is
$O(X\log^{-2\kappa-\kappa^2/32}X)+O(\sqrt X\log^2X)$.
\end{lemma}
\begin{proof}
Take $\tau=1-\kappa/2$, $u=1+\kappa/(4(1-\kappa))$, and $a=u-1$.
Markov's inequality gives relative exponent $E=(1-\kappa)(u-1)-\tau\log u$.
Since $\log(1+a)\ge a-a^2/2$,
\begin{equation}\label{eq:chernoff-half}
 E\le-\frac{a\kappa}2+\frac{\tau a^2}2
 =\frac{\kappa^2(-3+(7/2)\kappa)}{32(1-\kappa)^2}
 \le-\kappa^2/32.
\end{equation}
The factor $u^{-\tau\log\log X}$ multiplying the remainder is at most one.
\end{proof}

\begin{proof}[Proof of Theorem~\ref{thm:omega-consumer} and Remark~\ref{rem:tradeoff}]
Put $\rho=1/3750$, $z=\exp((\log X)^\rho)$.
For $\kappa=\rho$ use the hard weight of Theorem~\ref{thm:rough} and \eqref{eq:hard-positive-moment}.
For $0<\kappa<\rho$ fix $r=1-\kappa/\rho$ and use the weight of Theorem~\ref{thm:soft} and
\eqref{eq:soft-positive-moment}. Lemma~\ref{lem:soft-density} gives
$V^{\sfw}_{r,h}(z)\asymp_h W_r^2\asymp_h(\log X)^{-2\kappa}$ for this entire fixed-$r$ range, also when
$\kappa>\rho/2$ and hence $r<1/2$.
In either case each sign pattern has weighted mass comparable to $X\log^{-2\kappa}X$, and
\begin{equation}\label{eq:tradeoff-moment}
 \sum_{n\le X}\text{pair weight}\ u^{\omega(n)}
 \le CX(\log X)^{-2\kappa+(1-\kappa)(u-1)}+O(\sqrt X\log^2X).
\end{equation}
Here $\delta\rho=\kappa$ in the soft case.
Write $m=1-\kappa$, $\tau=m+\epsilon<1$, $a=\epsilon/\tau\in(0,1)$ and $u=1+a$.
The relative Markov exponent for $\omega(n)>\tau\log\log X$ is
\begin{equation}\label{eq:tradeoff-tail}
 ma-\tau\log(1+a)\le-\epsilon a+\tau a^2/2=-\epsilon^2/(2\tau)<0.
\end{equation}
The other endpoint has the same bound. Both unsigned tails are negligible relative to each sign-pattern mass.
The weight already excludes square factors at $p\le z$.
Failure of squarefreeness at larger primes costs at most
\begin{equation}\label{eq:large-square}
 (2X+h)\sum_{m>z}m^{-2}=O_h(X/z),
\end{equation}
using complete prefixes to count square multiples. Remove these exceptions as well. Then $\Omega=\omega$ at
both endpoints, and every remaining weight is at most one, so its cardinality is at least its weighted mass.
This proves \eqref{eq:tradeoff}. At the hard endpoint retain both roughness conditions and take
\begin{equation}\label{eq:tradeoff-numerical}
 \kappa=1/3750,\quad\epsilon=1/60000,\quad
 \kappa-\epsilon=1/4000,\quad 2\kappa=1/1875.
\end{equation}
This gives Theorem~\ref{thm:omega-consumer}.
\end{proof}

\begin{corollary}\label{cor:half-consumer}
At the maximal $z$ of Theorem~\ref{thm:soft}, every sign pair occurs for at least
$c_{h,r}X/(\log X)^{2\kappa}$ squarefree pairs with
\begin{equation}\label{eq:soft-half-consumer}
 \Omega(n),\Omega(n+h)\le(1-\kappa/2)\log\log X,\qquad\kappa=(1-r)/3750.
\end{equation}
At the hard endpoint the same argument gives $\Omega(n),\Omega(n+h)\le(1-1/7500)\log\log X$;
Theorem~\ref{thm:omega-consumer} gives the stated smaller upper bound.
\end{corollary}
\begin{proof}
Apply Lemma~\ref{lem:chernoff}, \eqref{eq:large-square} and the preceding pattern masses, or take
$\epsilon=\kappa/2$ in Remark~\ref{rem:tradeoff}.
\end{proof}

\section{The full prime-factor weight: phases, norms and mass}\label{sec:full}
In this and the remaining sections use padding $a(u)=4^{\omega(u)}$, $w=5^{\omega_{\Qset}}$,
$\Qset=\{p:y<p\le e^L,\ p\not\equiv1\pmod5,\ p\nmid h\}$, and the original padding probability
\begin{equation}\label{eq:original-rho}
 \rho_j(n,d)=\frac1{w(n)}\sum_{\substack{u\mid n,\ u\in\Qset^{\sfw}\\j\eta\le\log(du)<(j+1)\eta}}4^{\omega(u)}.
\end{equation}
The cutoff is $\rho_j\le K/L$ and the low-$\Qset$ cutoff is $400\log L$.
The original matrix has entry $L4^{\omega(u)}c_d(n)\ind_{u\mid n}\tau_j(n,d)\tau_j(m,d)/\sqrt{w(n)w(m)}$,
where now $\tau_j$ contains these two cuts and $\sigma_j$, and $m=n+hdu$.
No full $\Omega$ value enters this matrix or its finite-residue comparison.

\begin{proposition}[Numerical signs and two layers]\label{prop:phase}
Multiply the directed increasing edge with label $q'=du$ by any prescribed $\epsilon(q')\in\{\pm1\}$, and form
\begin{equation}\label{eq:phase-lift}
 \mathcal A_d=\begin{pmatrix}0&A_d^+\\(A_d^+)^*&0\end{pmatrix}.
\end{equation}
Duplicate $E$ in both layers. The moment and spectral estimates of
\cite[Theorem~5.4 and Proposition~8.3, pp.~24, 41]{OpenAI007} have the same permissible absolute constants,
with $K$ as row constant. Their finite comparison holds with $L^{18}$ by Proposition~\ref{prop:encoding}.
\end{proposition}
\begin{proof}
Check the original proof in order. The Hilbert--Schmidt expansion [O, pp.~24--25] projects each lifted closed
word to a physical word already counted there; it retains unequal consecutive tuples within each half.
The projected lifted words form a subset of the physical closed words, and each physical word has at most two
lifts, according to its starting layer. For a fixed numerical word its product of edge signs is a scalar of
modulus one, independent of residue draws. Remove it from the signed expectation before taking the absolute
value. Writing $X_w$ for the unphased contribution of a physical word, the positive word sums therefore obey
\begin{equation}\label{eq:phase-word-domination}
 \sum_{\widetilde w\in\mathcal W_{\rm lift}}|\E_*X_{\widetilde w}|
 \le2\sum_{w\in\mathcal W_{\rm phys}}|\E_*X_w|.
\end{equation}
This compares absolute word sums, as used in [O, p.~24], rather than signed operator norms.
The exact singleton identity [O, Lemma~5.5, p.~26] and the unlit, rank and singleton majorants
[O, pp.~26--33] are unchanged. The eligible numerical ranges, reciprocal weights and rare witnesses are the
ones verified in Proposition~\ref{lem:rare}; removing primes up to $y$ only prunes their positive sums.
The surviving perfect blocks and forest universe [O, Lemmas~7.1--7.2, pp.~34--37] are unchanged subsets.
The padding recursion [O, Lemma~7.3, p.~37] has oriented row bound $K$, with no numerical sign in the positive
row. Completion [O, pp.~37--38] now has external factor $2|\Dset|^2M\le e^{108L}$, the envelope already used
there, so the same absolute $C_2$ may be fixed. The signed finite expansion is covered by
Proposition~\ref{prop:encoding} before absolute values.
The arithmetic and counting proofs [O, Lemmas~5.2--5.3, 5.5--5.6, 6.1--6.5 and 7.1--7.2, pp.~22--37] use no
numerical value of $K$, $\eta$ or $A$. They use the bin eligibility only through $du\le e^{101L}$ and the degree
bound on $u$; these hold for $\eta\le1$. Within the padding-summation step, $K$ enters through the row recursion
of Lemma~7.3; its bound is then carried through the completion. The exponent $A$ belongs to the finite-origin
comparison, replaced here by Proposition~\ref{prop:encoding}.

For the one-orientation block $A_d^+$, weighted absolute row and column sums, with $\phi=\sqrt w$, satisfy
\begin{equation}\label{eq:phase-one-orientation-schur}
 \sum_m|A_d^+(n,m)|\frac{\phi(m)}{\phi(n)}\le Ka_d(n),\qquad
 \sum_n|A_d^+(n,m)|\frac{\phi(n)}{\phi(m)}\le Ka_d(m).
\end{equation}
Each is at most the corresponding two-orientation bound of [O, (8.4), p.~40].
Since $a_d$ is constant along an edge, the weighted Cauchy--Schwarz calculation is
\begin{equation}\label{eq:phase-schur-calculation}
\begin{aligned}
 \|A_d^+v\|^2
 &\le\sum_n Ka_d(n)\sum_m|A_d^+(n,m)|\frac{\phi(n)}{\phi(m)}|v_m|^2\\
 &\le K^2\sum_m a_d(m)^2|v_m|^2.
\end{aligned}
\end{equation}
Its transpose has the same bound. These are the Schur bounds on the invariant level sets of $a_d$;
the two-orientation proof in [O, Lemma~8.2, p.~40] uses their sum.
For $v=(v_1,v_2)$, with the duplicated projection $v=Ev$, this gives
\begin{equation}\label{eq:phase-schur}
 \|\mathcal A_d\|\le K\le2K,\qquad
 \sum_d\|\mathcal A_dv\|^2\le K^2(8W)^J\|v\|^2\le4K^2(8W)^J\|v\|^2.
\end{equation}
Thus the original permissible two-orientation constants still apply.
If $\mathcal A=\sum_d\mathcal A_d$ and $A^+=\sum_d A_d^+$, the exact testing identity for real vectors is
\begin{equation}\label{eq:phase-quadratic-form}
 \langle(g_1,g_2),\mathcal A(g_1,g_2)\rangle
 =\langle g_1,A^+g_2\rangle+\langle g_2,(A^+)^*g_1\rangle
 =2\langle g_1,A^+g_2\rangle.
\end{equation}
All edge signs and testing vectors here are real. In particular this recovers twice the ordered increasing-edge
bilinear form for the three sign-exponent cases.
Input~\ref{input:transfer} and the Markov proof of [O, Proposition~8.3, p.~41] give
$K(C_3\sqrt W)^J$ outside $e^{-2k}$ of origins, with
$C_3=3\max(2,2\sqrt8,eC_2)$ as in [O, (8.7)]. This proves the stated constants without comparing signed
submatrix norms by entrywise domination.
\end{proof}

\begin{proposition}[Positive actual-integer norms]\label{prop:full-norm}
Uniformly for $99/100\le r\le1$, $2\le y\le e^{L^{1/12}}$,
$e^{L^{18}/2}\le Y\le e^{2L^{18}}$, put $B_r(n)=r^{\Omega(n)}S_y(n)$ and
$S_r=\prod_{p\in\Qset}(1+4r^2/p)$. Then
\begin{align}
 \sum_{n\le Y}B_r(n)^2w(n)&\ll YS_r(\log Y)^{r^2-1},&
 \sum_{n\le Y}B_r(n)^2&\ll Y(\log Y)^{r^2-1},\label{eq:full-norm}\\
 \sum_{n\le Y}B_r(n)^2e^{\omega_{\mathcal T}(n)}
 &\ll Y(\log Y)^{r^2-1}
 \exp\left((e-1)r^2\sum_{p\in\mathcal T}1/p\right)&&\label{eq:full-tilted-norm}
\end{align}
for any prime set $\mathcal T\subset[2,e^L]$.
\end{proposition}
\begin{proof}
Shiu's nonnegative mean theorem in the form reproduced in \cite[p.~1]{Henriot} bounds a mean by
$CY\exp(\sum_{p\le Y}(F(p)-1)/p)$ for functions of fixed divisor growth.
For the two functions in \eqref{eq:full-norm}, $F\le5^{\omega(n)}\le\tau_5(n)$; for
\eqref{eq:full-tilted-norm}, $F\le3^{\omega(n)}\le\tau_3(n)$.
Thus $F(p^j)\le A^j$ and $F(n)\le B_\xi n^\xi$ have uniform constants for every fixed $\xi>0$.
Apply the theorem with $q=1$, interval and progression exponents $1/2$ on dyadic long blocks above $\sqrt Y$.
The small prefix costs $O(\sqrt Y\log^4Y)$.
The prime values are $5r^2$ on $\Qset$, $r^2$ elsewhere, or $er^2$ on $\mathcal T$ and $r^2$ elsewhere.
The squarefree cutoff changes only higher powers. Uniformly summable $O(p^{-2})$ errors identify the extra
$\Qset$ product with $S_r$, and Mertens gives the displayed powers of $\log Y$.
These are bounds on actual integers, rather than a residue model for full $\Omega$.
\end{proof}

\begin{proposition}[Unsigned full-weight pair mass]\label{prop:full-mass}
For fixed even $h$, uniformly for $99/100\le r\le1$ and
$2\le y\le\exp((\log X)^{1/216})$,
\begin{equation}\label{eq:full-pair-mass}
 c_hX(\log X)^{2r-2}\le\sum_{n\le X}r^{\Omega(n)+\Omega(n+h)}S_y(n)S_y(n+h)
 \le C_hX(\log X)^{2r-2}
\end{equation}
for sufficiently large $X$.
\end{proposition}
\begin{proof}
Choose a cutoff $z_0$ and independently for each $p\le z_0$ choose two thresholds with
$\Prob(T_i>k)=r^k$, $k\ge0$; at $r=1$ take $T_i=\infty$.
For $p\le y$ replace each threshold by its minimum with $2$.
The complement of the event $p^{T_1}\mid n$ or $p^{T_2}\mid n+h$ has expectation exactly the two local
prime-power weights, including the squarefree cap. This is a nonnegative sieve measure on integers and
threshold choices, within the probability setup of \cite[pp.~1--3]{Ford}.

For squarefree $d$, expand each bad-prime union into at most three conjunction terms.
Conditional on thresholds a consistent term is one prime-power CRT progression, counted with error $O(1)$
regardless of its modulus; inconsistent terms are zero. Average to obtain
\begin{equation}\label{eq:threshold-crt}
 A_d=Xg(d)+O(3^{\omega(d)}),\qquad g(d)=\prod_{p\mid d}g(p).
\end{equation}
Infinite thresholds may be truncated with an infinity option; bounded convergence then applies to each finite
sieve sum. The divisor upper/lower identities of Input~\ref{input:fl} hold for every sample of the bad-prime
set and may be integrated.

For $p\nmid h$, $p>y$, a $p$-adic count gives
\begin{equation}\label{eq:full-uncapped-local}
 1-g(p)=1-\frac{2(1-r)}{p-r}.
\end{equation}
Indeed at one of the two distinct roots the valuation-$k$ density is $(1-1/p)p^{-k}$, and summing
$1-2/p+2(1-1/p)\sum_{k\ge1}r^kp^{-k}$ gives the formula.
For $p\le y$ the factor is \eqref{eq:soft-local}. The finitely many $p\mid h$ have positive local expectations.
All bad events are subsets of the ordinary two-root event, so $g(p)\le\nu_h(p)/p<1$.
In the uncapped case one may also use $1-r^k\le(1-r)k$ and the mean valuation $2/(p-1)$.
Thus the sieve dimension condition is uniformly bounded by dimension two. Logarithms of generic factors have
coefficient $-2(1-r)/p$; the squarefree cap changes them by $O(p^{-2})$, and fixed exceptional factors stay
uniformly positive for $r\ge.99$. Hence
\begin{equation}\label{eq:full-local-product}
 \prod_{p\le z_0}(1-g(p))\asymp_h(\log z_0)^{2r-2}\qquad(y\le z_0).
\end{equation}
Take a fixed sufficiently large $s_0$ with lower-sieve relative error at most $1/4$,
$D=X^{1/4}$ and $z_0=X^{1/(4s_0)}$. Eventually $y\le z_0\le D^{1/2}$.
The total CRT error is $O(D\log^2D)$, smaller than the main term.
The finite prime-weighted mass is thus bounded below by a positive multiple of $X(\log X)^{2r-2}$.
At each endpoint at most $4s_0+1$ prime factors counted with multiplicity exceed $z_0$.
Restoring them costs at most the fixed factor
$r^{8s_0+2}\ge(.99)^{8s_0+2}>0$, proving the lower bound.

For the upper bound drop $S_y$ and apply \cite[Theorem~1, pp.~1--2]{Henriot} to the same fixed polynomials
$t,t+h$ and $F(a,b)=r^{\Omega(a)+\Omega(b)}$.
All function-class bounds hold uniformly since $0\le F\le1$; the polynomial, interval and coefficient
hypotheses were verified in Proposition~\ref{prop:positive-moment}.
The positive Euler factors give coefficient $(2r-2)/p$ with uniformly summable higher-power errors.
Dyadic summation completes the bound. The restoring constant is not asserted uniform as $r$ tends to zero.
\end{proof}

\section{External padding collision}\label{sec:full-collision}

\begin{inputthm}[Hal\'asz in the Montgomery--Tenenbaum form, {\cite[Chapter~III.4]{Tenenbaum}}]\label{input:halasz}
For $Y\ge3$ and $1\le T\le Y/2$, if $F$ is multiplicative, $F(1)=1$ and $|F(n)|\le1$, then
\begin{equation}\label{eq:halasz}
 \frac1Y\left|\sum_{m\le Y}F(m)\right|
 \ll(1+M_F)e^{-M_F}+T^{-1/2},\quad
 M_F=\min_{|v|\le2T}\sum_{p\le Y}\frac{1-\Re(F(p)p^{-iv})}p.
\end{equation}
A rescaling of the height parameter gives the convention $|v|\le2T$ in this display.
Our use has $T=L^{1000}$ and $Y\ge e^{L^{18}/2}$, so the stated height and length requirements hold for large $L$.
\end{inputthm}

\begin{proposition}[Common-prefix collision]\label{prop:external-collision}
Use the range and notation of Proposition~\ref{prop:full-norm}. Conditional on an integer $um$, choose a
squarefree $\Qset$ divisor $v$ with probability $4^{\omega(v)}/w(um)$. Then
\begin{equation}\label{eq:external-collision}
 \sum_{u\in\Qset^{\sfw}}\frac{4^{\omega(u)}r^{2\omega(u)}}u
 \sum_{m\le Y}B_r(m)^2\Prob_{v\mid um}(|\log(u/v)|\le1)
 \ll_h YS_r(\log Y)^{r^2-1}\frac{\log(2L)}L.
\end{equation}
The $u$-sum is finite and has the same physical prefix $m\le Y$ for every $u$.
\end{proposition}
\begin{proof}
Replace the probability by $\E_v\exp(it\log(u/v))$ and call the resulting sum $H(t)$, $|t|\le2$.
For a $\Qset$ prime put $z_p=p^{it}$, $k_p=4r^2/p$ and
\begin{equation}\label{eq:external-local-factors}
 C_{0,p}=1+\frac{k_p(4+z_p)}5,\qquad
 C_{1,p}=\frac{1+4z_p^{-1}+k_p(z_p+4)}5.
\end{equation}
If $p\nmid m$, not choosing it in $u$ contributes $1$; choosing it contributes
$k_p(z_p/5+4/5)$, since the conditional $v$ probabilities are $1/5,4/5$.
If $p\mid m$, these same $v$ probabilities apply whether or not it is chosen in $u$, giving $C_{1,p}$.
The finite Euler identity is therefore
\begin{equation}\label{eq:external-euler}
 H(t)=\left(\prod_{p\in\Qset}C_{0,p}\right)\sum_{m\le Y}F_t(m),\qquad
 F_t(p^j)=\begin{cases}B_r(p^j)^2C_{1,p}/C_{0,p},&p\in\Qset,\\B_r(p^j)^2,&p\notin\Qset.
 \end{cases}
\end{equation}
Define $F_t$ multiplicatively with $F_t(1)=1$.
For $c=\Re z_p$, direct algebra gives
\begin{equation}\label{eq:external-abs-identity}
 |C_{0,p}|^2-|C_{1,p}|^2=\frac8{25}(1-c)(1+5k_p+2k_pc)\ge0,\qquad \Re C_{0,p}\ge1.
\end{equation}
Thus $|F_t|\le1$, as required in Input~\ref{input:halasz}. Its prime values are
\begin{equation}\label{eq:external-prime-values}
 F_t(p)=\tfrac{r^2}5(1+4p^{-it})+O(p^{-1})\ (p\in\Qset),\qquad F_t(p)=r^2\ (p\notin\Qset).
\end{equation}
The harmonic sum of these coefficient errors is $O(1)$.

Put $v_0=\log(2y)$ and $\Psi(t)=\log((1+L|t|)/(1+v_0|t|))$.
Fixed-modulus PNT and partial summation give, uniformly for $|t|\le3$,
\begin{align}
 \sum_{p\in\Qset}\frac{1-\cos(t\log p)}p&=\tfrac34\Psi(t)+O_h(1),\label{eq:external-harmonic}\\
 \left|\prod_{p\in\Qset}C_{0,p}\right|&\ll_h S_r\exp(-3r^2\Psi(t)/5).\label{eq:external-prefactor}
\end{align}
For the latter, the leading logarithmic difference from $1+4r^2/p$ is
$-(4r^2/5p)(1-\cos(t\log p))$; all $O(p^{-2})$ errors are summable.
This is the same partial summation as \cite[(4.12), p.~18]{OpenAI007}, now with the lower cutoff $y$.

Take $T=L^{1000}$. For any candidate twist $v$, \eqref{eq:external-prime-values} expands the distance as
\begin{equation}\label{eq:external-distance}
 (1-r^2)\log\log Y+r^2S_Y(v)-\tfrac45r^2S_Q(v)+\tfrac45r^2S_Q(v+t)+O_h(1),
\end{equation}
where $S_Y(v)=\sum_{p\le Y}(1-\cos(v\log p))/p$ and $S_Q$ is the $\Qset$ sum.
If $|v|\le1$, partial summation of PNT gives $S_Y(v)=\log(1+|v|\log Y)+O(1)$.
The increasing concave function $x\mapsto\log((1+Lx)/(1+v_0x))$ vanishes at zero, and is subadditive.
Thus $\Psi(t)\le\Psi(v)+\Psi(v+t)$. Since $\log Y\ge L^{18}/2$, for sufficiently large $L$,
\begin{equation}\label{eq:external-small-twist}
 \log(1+|v|\log Y)\ge\tfrac65\log(1+L|v|)\ge\tfrac65\Psi(v)\qquad(|v|\le1).
\end{equation}
One may verify the inequality separately for $L|v|\le1$ and $1\le L|v|\le L$.
Using \eqref{eq:external-harmonic} in \eqref{eq:external-distance} therefore gives
\begin{equation}\label{eq:external-distance-lower}
 M_{F_t}\ge(1-r^2)\log\log Y+\tfrac35r^2\Psi(t)-O_h(1)
\end{equation}
for these twists.
For $1\le|v|\le2T$, work on $\sigma=1+1/\log Y$. Replacing the sharp prime sum by its $p^{-1/\log Y}$
smoothing costs $O(1)$, by partial summation; higher powers cost $O(1)$ too.
The Euler product and Euler summation with cutoff comparable to $|v|+3$ give
$|\zeta(\sigma+iv)|\ll\log(|v|+3)$, hence
$S_Y(v)\ge\log\log Y-O(\log\log(3T))$.
Also $S_Q(v)\le2\sum_Q1/p\le(3/2)\log L+O_h(1)$, $S_Q(v+t)\ge0$, and
$\log\log Y\ge18\log L-O(1)$. These inequalities imply \eqref{eq:external-distance-lower} in the large-twist
range as well. This checks the full Hal\'asz minimum, not only a zero twist.

Apply \eqref{eq:halasz}, \eqref{eq:external-prefactor}, and
$M_{F_t}\le2\log\log Y+O(1)$ to obtain
\begin{equation}\label{eq:external-char-bound}
 |H(t)|\ll_h YS_r(\log Y)^{r^2-1}\log(2L)
 \left(\frac{1+v_0|t|}{1+L|t|}\right)^{6r^2/5}.
\end{equation}
The $T^{-1/2}=L^{-500}$ term is absorbed uniformly; the other density and phase factors have magnitude at least
$L^{-2}$ in the stated compact $r$ and $Y$ ranges.
Finally apply the positive kernel \eqref{eq:triangular} to $\log(u/v)$, interchange only finite sums, and use
\begin{equation}\label{eq:external-integral}
 \int_0^2\left(\frac{1+v_0t}{1+Lt}\right)^\nu dt
 \ll L^{-1}+(v_0/L)^\nu\ll L^{-1},\qquad
 \nu=6r^2/5>1,\quad v_0\le2L^{1/12},\quad(11/12)\nu>1.
\end{equation}
This proves \eqref{eq:external-collision}, retaining complex phases throughout the Euler identity and distance
minimization.
\end{proof}


\begin{lemma}[Uniform upper sieve for marked primes]\label{lem:marked-window}
There is an absolute $u_0>0$ such that, for an interval $I$ of $N$ consecutive positive integers, $z\ge2$ and
$\log N/\log z\ge u_0$, every set $\mathcal R\subset\{p\le z\}$ satisfies
\begin{equation}\label{eq:marked-window-sieve}
 \#\{n\in I:p\nmid n\text{ for every }p\in\mathcal R\}
 \le2N\prod_{p\in\mathcal R}(1-1/p)+O(N^{1/2}\log^2(2N)),
\end{equation}
with an absolute error constant. Consequently, for $0\le\delta\le1$, $r=1-\delta$,
\begin{equation}\label{eq:marked-window-average}
 \sum_{n\in I}r^{\omega_z(n)}
 \ll N(\log(2z))^{-\delta}+N^{1/2}\log^2(2N).
\end{equation}
\end{lemma}
\begin{proof}
Use Selberg's upper sieve with divisor level $U=N^{1/4}$.
Here is its finite quadratic construction, which also verifies uniformity in the marked set.
Let $\mathcal D_U=\{d\le U:d\mid\prod_{p\in\mathcal R}p\}$ and
\begin{equation}\label{eq:selberg-GF}
 G_{\mathcal R}(U)=\sum_{d\in\mathcal D_U}\frac1{\varphi(d)},\qquad
 F_{\mathcal R}=\prod_{p\in\mathcal R}(1-1/p)^{-1}.
\end{equation}
Put $y_q=\mu(q)/(\varphi(q)G_{\mathcal R}(U))$ for $q\in\mathcal D_U$, and define
$\lambda_d/d=\sum_{q\in\mathcal D_U,\ d\mid q}\mu(q/d)y_q$.
Finite divisor inversion gives $\lambda_1=1$ and $y_q=\sum_{d\in\mathcal D_U,\ q\mid d}\lambda_d/d$.
Since $(d,e)=\sum_{q\mid(d,e)}\varphi(q)$, the quadratic main term is exactly
\begin{equation}\label{eq:selberg-diagonalization}
 \sum_{d,e\in\mathcal D_U}\frac{\lambda_d\lambda_e}{[d,e]}
 =\sum_{q\in\mathcal D_U}\varphi(q)y_q^2=G_{\mathcal R}(U)^{-1}.
\end{equation}
Moreover
\begin{equation}\label{eq:selberg-weight-bound}
 \lambda_d=\frac{\mu(d)d}{\varphi(d)G_{\mathcal R}(U)}
 \sum_{\substack{m\le U/d\\m\mid\prod_{p\in\mathcal R}p,\ (m,d)=1}}\frac1{\varphi(m)},\qquad
 |\lambda_d|\le d/\varphi(d)\le\tau(d).
\end{equation}
The sifted indicator is bounded by $(\sum_{d\mid n}\lambda_d)^2$.
Counting each modulus $[d,e]$ in $I$ as $N/[d,e]+O(1)$ therefore gives
$N/G_{\mathcal R}(U)+O(U^2\log^2(2U))$, independently of the interval position and marked set.

For the Rankin tail take $\sigma=1/\log z$. The full squarefree divisor sum is $F_{\mathcal R}$, and
\begin{equation}\label{eq:selberg-rankin-tail}
\begin{aligned}
 \frac{F_{\mathcal R}-G_{\mathcal R}(U)}{F_{\mathcal R}}
 &\le U^{-\sigma}\prod_{p\in\mathcal R}\left(1+\frac{p^\sigma-1}{p}\right)\\
 &\le\exp\left(-\frac{\log U}{\log z}+C_0\right).
\end{aligned}
\end{equation}
The constant $C_0$ is absolute: $p^\sigma-1\le(e-1)\log p/\log z$ for $p\le z$, and
$\sum_{p\le z}(\log p)/p\ll\log z$ by Mertens' theorem \cite[Section~2.2]{MV}.
Choose $u_0\ge4(C_0+\log2)$ once and for all. Then $G_{\mathcal R}(U)\ge F_{\mathcal R}/2$, proving
\eqref{eq:marked-window-sieve}. This proves the sieve bound for every marked set before any averaging.

Independently mark each prime up to $z$ with probability $\delta$.
At a fixed $n$, the probability that no marked prime divides it is $r^{\omega_z(n)}$.
Average \eqref{eq:marked-window-sieve}; its error is uniform and
\begin{equation}\label{eq:marked-product-expectation}
 \E_{\mathcal R}\prod_{p\in\mathcal R}(1-1/p)
 =\prod_{p\le z}(1-\delta/p)\asymp(\log(2z))^{-\delta}.
\end{equation}
The comparison follows from the logarithmic expansion and Mertens, uniformly in $0\le\delta\le1$.
This gives \eqref{eq:marked-window-average}.
\end{proof}

\section{Cancellation with the full weight}\label{sec:full-proof}

\begin{proof}[Proof of Theorem~\ref{thm:full}]
In this proof abbreviate $C_\Omega,W_\Omega,\alpha_\Omega,\delta_\Omega$ by $C,W,\alpha,\delta$.
Thus \eqref{eq:full-fixed} fixes $W=16eC^2$, $\alpha=1/(1200W)$, $\delta=\sqrt{\alpha/216}$ and $r=1-\delta$.
Set
\begin{equation}\label{eq:full-parameters}
\begin{gathered}
 L=(\log X)^{1/18},\quad J=\lfloor\alpha\log L\rfloor,\quad\eps=\alpha/100,\\
 b=1+r^{-2},\quad E_0=18\delta^2=\alpha/12,\quad
 \eta=L^{-\eps},\quad K=L^{2\alpha\log b+2E_0+4\eps}.
\end{gathered}
\end{equation}
The exponent of $K$ is less than $4\alpha$; these are fixed powers of $L$, as required in
Proposition~\ref{prop:encoding}. The centre supply is Definition~\ref{def:graph}, with the same $H_0$ and
$W\le V_i\le2W$. All label primes exceed $y$.

Put $f_i(n)=\lambda(n)^{e_i}B(n)$. For $q'=du$, complete additivity of $\Omega$ and complete multiplicativity of
$\lambda$ give the exact identity
\begin{equation}\label{eq:full-dilation}
 f_1(q'm)f_2(q'(m+h))
 =\lambda(q')^{e_1+e_2}r^{2\Omega(q')}f_1(m)f_2(m+h).
\end{equation}
The small-prime-square mask is unchanged, with no coprimality condition on $q'$ and $m$.
Multiply the edge by $\lambda(q')^{e_1+e_2}$ and use Proposition~\ref{prop:phase}.
The raw term has positive harmonic mass
\begin{equation}\label{eq:full-raw-mass}
 S_{0,r}=\sum_j\sum_{(d,u)\in\mathcal S_j}\frac{4^{\omega(u)}r^{2J+2\omega(u)}}{du},\qquad
 \tfrac12r^{2J}S_rV_*\le S_{0,r}\le r^{2J}S_rV_*.
\end{equation}
To prove the lower bound, its padding reciprocal law selects a prime with probability
$4r^2/(p+4r^2)$, with expected logarithmic size at most $4L$ and degree at most $4\log L$.
Use the two Markov exclusions in Proposition~\ref{lem:mass}, now with this law.
The raw cofactor cutoff remains in $(Xe^{-\eta},X]$.
The Nair--Tenenbaum upper theorem gives mass
$O_h(\eta X(\log X)^{-2\delta})$ in the full missing interval of length comparable to $\eta X$.
Here the weights $r^{\Omega}$ are bounded by one, the polynomials are the fixed $t,t+h$, and
$\eta X>X^{1/2}$ eventually, so the interval and function-class hypotheses checked in
Proposition~\ref{prop:positive-moment} hold. The normalized raw error is
$O_h(\eta(\log X)^{-2\delta})$.

We next align all physical prefixes before enlarging padding sums.
The positive centred-factor majorant is
\begin{equation}\label{eq:full-majorant}
 a_d^+(n)=\prod_{p\mid d}(\ind_{p\mid n}+1/p)
 =\sum_{e\mid d}\frac ed\ind_{e\mid n}.
\end{equation}
In a term with $u\mid n$ and $e\mid n$, write $n=uem$; the prime supports of $u,e$ are disjoint and exceed $y$.
Then $B(uem)^2=r^{2\omega(u)+2\omega(e)}B(m)^2$, and the bin gives
\begin{equation}\label{eq:alignment}
 Xe^{-\eta}<T_j/(ud)\le X,\qquad T_j^{-1}\le e^\eta/(Xud),\qquad m\le Xd/e.
\end{equation}
Its coefficient in the normalized one-endpoint squared edge mass is bounded by
$e^\eta e r^{2\omega(e)}/(Xd^2)$ times the external sum
$\sum_u4^{\omega(u)}r^{2\omega(u)}/u$ and the $B(m)^2$ sum through $Y=Xd/e$.
Only now extend $u$ to all of $\Qset^{\sfw}$. Apply \eqref{eq:full-norm} without $w$, cancel the prefix $Y$
against this coefficient and sum $d,e$. Since $\sum_{e\mid d}r^{2\omega(e)}=(1+r^2)^J$, the whole squared mass is
\begin{equation}\label{eq:full-positive-edge}
 \ll_h S_r(\log X)^{r^2-1}(1+r^2)^JV_*.
\end{equation}
Every eligible $(d,u)$ belongs to one bin, so this alignment introduces no bin multiplicity.
At the arrival, $n'=n+hdu=ue(m+hd/e)$ and $e\mid d$.
Its cofactor $m'=m+hd/e$ is at most $(X+h)d/e$. Enlarge to this common prefix and repeat the same argument.
This proves \eqref{eq:full-positive-edge} at either endpoint, with logarithms comparable to $\log X$.

For a padding-cut failure, apply Cauchy--Schwarz in the positive edge measure.
One squared factor has the bad indicator and $B(n)^2$, the other discards the indicator and uses
\eqref{eq:full-positive-edge}. The original $\rho_j$ reads only the active $\Qset$ primes, so
$\rho_j(uem,d)=\rho_j(um,d)$. Since the eligible $u$ and every $v$ in its bin have
$|\log(u/v)|\le\eta\le1$,
\begin{equation}\label{eq:full-bad-padding}
 \ind_{\rho_j(um,d)>K/L}\le(L/K)\rho_j(um,d)
 \le(L/K)\Prob_{v\mid um}(|\log(u/v)|\le1).
\end{equation}
After \eqref{eq:alignment} the prefix is the same $Y=Xd/e$ for all $u$.
Proposition~\ref{prop:external-collision} makes the squared bad factor smaller than
\eqref{eq:full-positive-edge} by $O_h(\log(2L)/K)$.
At the arrival use $m'=m+hd/e$ and the common prefix just proved; the same inequality applies.
Dividing the resulting two-endpoint deletion by \eqref{eq:full-raw-mass} gives
\begin{equation}\label{eq:full-padding-deletion}
 \ll_h(\log X)^{r^2-1}b^J\sqrt{\log(2L)/K}.
\end{equation}

For a $\Qset$-degree failure, retain $\omega(u)\le100\log L$ until
$\omega_{\Qset}(uem)>400\log L$ implies $\omega_{\Qset}(m)>300\log L$.
Markov and \eqref{eq:full-tilted-norm} give a squared-factor saving $O(L^{-200})$, since
$\sum_Q1/p\le\log L+O_h(1)$; only afterward enlarge the $u$-sum.
Cauchy--Schwarz saves $L^{-100}$.
For a $\Pset$-degree failure, $\omega_{\Pset}(uem)\le J+\omega_{\Pset}(m)$ and
\begin{equation}\label{eq:full-p-tail}
 e^{-(6W-1)J}\exp(2(e-1)WJ)\le e^{-2WJ}\qquad(W\ge10).
\end{equation}
This gives $e^{-WJ}$ after Cauchy--Schwarz.
For rare-site failures discard $B(n)B(n+hdu)\le1$ and use the positive deletion proof of
\cite[Lemma~4.6, pp.~19--20]{OpenAI007}, with \eqref{eq:rare} and Proposition~\ref{prop:encoding}.
The ratio of $SV_*$, where $S=\prod_{p\in\Qset}(1+4/p)$, to $S_{0,r}$ and the inverse mass density are fixed powers of $L$, absorbed in its
exponential error. The complete relative deletion error is
\begin{equation}\label{eq:full-deletion-budget}
 \ll_h L^{E_0}b^J\{\sqrt{\log(2L)/K}+L^{-100}+e^{-WJ}\}+e^{-L^{0.8}}.
\end{equation}
Here $(\log X)^{(r^2-1)+2\delta}=(\log X)^{\delta^2}=L^{E_0}$.
The nonrare estimates were actual-integer positive norms, not a finite-residue model for $B$.

For the spectral term use Proposition~\ref{prop:phase} and the unweighted signed-trace comparison of
Proposition~\ref{prop:encoding}.
The actual squared testing norms of $\sqrt w f_i$ average to
$O_h(MS_r(\log X)^{r^2-1})$, by Proposition~\ref{prop:full-norm}; Cauchy--Schwarz handles the bilinear form.
The overlap identity \eqref{eq:overlap} and its deterministic exceptional-block bound apply to the two layers.
All normalizations are fixed powers of $L$, so their exponential errors remain negligible.
After the bins and \eqref{eq:full-raw-mass}, the relative spectral term is
\begin{equation}\label{eq:full-spectral}
 \ll_h L^{E_0}\frac K\eta\left(\frac C{r^2\sqrt W}\right)^J+e^{-L^{0.8}}.
\end{equation}
The $S_r$ in the actual norm cancels the $S_r$ in the raw mass.

Finally repeat the centering proof of \cite[Lemmas~4.1--4.2, 4.4, pp.~13--16]{OpenAI007} with the present $f_i$.
At least one signed factor has prime value $-r$ at every prime, since $S_y$ changes only higher powers.
For that factor,
\begin{equation}\label{eq:full-signed-distance}
 M(f_i;Y,Q)=rM(\lambda;Y,Q)+(1-r)\sum_{p\le Y}1/p
 \ge\tfrac14\log\log Y-O(1).
\end{equation}
Use precisely Input~\ref{input:mrt}; $Y\ge e^{L^{18}/2}$, $H_0/2\le D\le e^{2L}\ll Y$, and
$Q=\min((\log Y)^{1/125},(\log D)^5)$. The height is $|t|\le Y$.
Its errors give the short Fourier bound $DL^{-\Gamma}$ with
\begin{equation}\label{eq:full-centering-parameters}
 a=994999/1000000,\qquad\Gamma=2571/100000,\qquad s=1/1000000,
 \qquad a<199/200,\quad\Gamma<18/700.
\end{equation}
In choosing the final threshold also require
\begin{equation}\label{eq:marked-window-scale}
 \tfrac12L^{199/200-a}\ge u_0,\qquad z_{\rm m}=e^{L^a},\qquad
 \frac{\log D}{\log z_{\rm m}}\ge u_0
 \quad(H_0/2\le D\le e^{2L}),
\end{equation}
where $u_0$ is fixed in Lemma~\ref{lem:marked-window}. This is possible because $199/200-a=10^{-6}>0$.
The distance is eventually at least $(18/5)\log L$. The short error is
$O(L^{-199/200}\log L)$, and the long error is $O(L^{-18/700})$.
Proposition~\ref{lem:shift-polynomial} has $\|\mathcal B\|_\infty\ll L^{-a}$ and fourth moment
$O(D^{-1}L^{-4a})$ at this changed cutoff.

The averaged squared short-window norms are bounded by $D(\log Y)^{r^2-1}$:
count each integer in at most $D$ windows and apply \eqref{eq:full-norm} without $w$, followed by Cauchy--Schwarz
in the origin average. For the pointwise other factor, use
\begin{equation}\label{eq:full-pointwise}
 B(m)\le r^{\omega_{p\le e^{L^a}}(m)},\qquad |G_v|\ll_h DL^{-a\delta}.
\end{equation}
For the pointwise assertion, take $N=(2h+1)D$ and $z=z_{\rm m}$ in Lemma~\ref{lem:marked-window}.
Condition \eqref{eq:marked-window-scale} gives $\log N/\log z_{\rm m}\ge u_0$.
For every marked-prime set the Selberg sieve has divisor level $N^{1/4}$ and the Rankin estimate gives
$G_{\mathcal R}(N^{1/4})\ge F_{\mathcal R}/2$; average only after applying this uniform bound.
The expected product is $\prod_{p\le z_{\rm m}}(1-\delta/p)\asymp L^{-a\delta}$.
The remainder $O_h(D^{1/2}\log^2(2D))$ is smaller than $DL^{-a\delta}$ throughout the stated $D$ range,
for sufficiently large $L$. Dropping the squarefree mask and taking absolute values of the endpoint function
therefore proves $|G_v|\ll_h DL^{-a\delta}$ in every positive window. The separately required CRT estimate in
Proposition~\ref{lem:shift-polynomial} remains $D^{-1}e^{O(L^a\log L)}$, negligible because $a<199/200$.

For a nonempty partial-centred subset $I$, put $q_I=\prod_{i\in I}p_i$; the physical shift is $hq_I$.
The subset extracts only $J-|I|$ centre primes.
Relative to $r^{2J}$ in the raw mass its missing factor is $r^{-2|I|}$, so summing subsets gives
$(1+r^{-2})^J=b^J$. All remaining numerical signs have modulus one and leave the shift-polynomial estimate
unchanged. Padding extraction has harmonic mass $S_r$.
The Fourier split $|\mathcal B|\le L^{-1-s}$ now gives the two relative terms
\begin{equation}\label{eq:full-centering-budget}
 L^{E_0+\eps+\alpha\log b-s},\qquad
 L^{5(1-a)+4s-\Gamma+(36-a)\delta+\eps+\alpha\log b}.
\end{equation}
The squared-norm loss is $L^{E_0}$; the pointwise loss is $L^{36\delta-a\delta}$.
If only the second factor is signed, reindex by $hq_I$ before Fourier averaging, put the signed factor in $F$, and
allow negative offsets in $G$. Removing the initial $O_h(D)$ origins makes all arguments positive and costs
$O_h(D/Y)$, exponentially little. Hence this centering bound covers all three nonzero exponent pairs.
The input is the classical MRT distance bound; the half-plane $\Re s>7/8$ is not used.

We record the numerical margins. Since $C_3\ge6\sqrt8$, $C^2\ge288$; using $e>8/3$ gives
$\alpha<1/14700000$, $\delta<1/56000$ and $r\ge.99$.
Also $\log b<.71$ and
\begin{equation}\label{eq:full-spectral-margin}
 \log(4\sqrt e\,r^2)-2\log b
 =\tfrac12+6\log r-2\log((1+r^2)/2)>2/5.
\end{equation}
For example $\log r\ge-1/99$ and the last logarithm is nonpositive.
The exponent in \eqref{eq:full-spectral} is at most
\begin{equation}\label{eq:full-spectral-exponent}
 3E_0+5\eps-(2/5)\alpha=\alpha/4+\alpha/20-2\alpha/5=-\alpha/10.
\end{equation}
The padding deletion exponent in \eqref{eq:full-deletion-budget} is $-2\eps$ before its square-root logarithm;
the raw error has exponent $-\eps$, and both degree terms are smaller.
For the two centering terms, $E_0+2\eps+\alpha\log b<s$, and
\begin{equation}\label{eq:full-centering-margin}
 5(1-a)+4s+36/56000+1/10000000<.02571=\Gamma;
 \qquad\text{the left side is less than } .025652.
\end{equation}
Thus both exponents in \eqref{eq:full-centering-budget} are less than $-\eps$.
Absorb fixed logarithmic factors into $\eps/2$ to obtain
\begin{equation}\label{eq:full-final-saving}
 \left|\sum_{n\le X}f_1(n)f_2(n+h)\right|
 \ll_h X(\log X)^{-2\delta}L^{-\eps/2},\qquad
 L^{-\eps/2}=(\log X)^{-\alpha/3600}.
\end{equation}
This is \eqref{eq:full-main}; Proposition~\ref{prop:full-mass} proves \eqref{eq:full-mass-main}.
All constants were fixed before $h,X$; only the final lower threshold depends on $h$.
\end{proof}

\begin{corollary}\label{cor:full-patterns}
With \eqref{eq:full-fixed}, let $M_0(X,y)=\sum_{n\le X}B(n)B(n+h)$.
For each sign pair the weighted mass is
\begin{equation}\label{eq:full-pattern-mass}
 \sum_{\substack{n\le X\\(\lambda(n),\lambda(n+h))=(s,t)}}B(n)B(n+h)
 =\tfrac14M_0(X,y)\bigl(1+O_h((\log X)^{-\gamma_\Omega})\bigr).
\end{equation}
Every sign pattern also occurs for $\gg_h X/(\log X)^{2\delta_\Omega}$ integers $n\le X$ with both endpoints
squarefree and
\begin{equation}\label{eq:full-omega-consumer}
 \Omega(n),\Omega(n+h)\le(1-\delta_\Omega/2)\log\log X.
\end{equation}
\end{corollary}
\begin{proof}
Expand \eqref{eq:sign-expansion}, use all three signed estimates in Theorem~\ref{thm:full}, and divide by
\eqref{eq:full-pair-mass}; no exact unsigned Euler leading constant is asserted.
For the cardinality assertion take $y=\exp((\log X)^{1/216})$, $\tau=1-\delta/2$ and $u=1+\delta/(4r)$.
Then $ru=1-3\delta/4<1$. The Nair--Tenenbaum theorem for the nonnegative bounded function
$(ru)^{\Omega(n)}r^{\Omega(n+h)}$ gives
$O_h(X(\log X)^{ru+r-2})$; its fixed-polynomial hypotheses are those already verified, and squarefree masks
may be dropped for this upper bound.
Relative to the unsigned mass the Markov exponent is $r(u-1)-\tau\log u$.
Equation~\eqref{eq:chernoff-half} with $\kappa=\delta$ makes it at most $-\delta^2/32$.
The second endpoint is identical. The squarefree support at small primes and \eqref{eq:large-square} with $z=y$
remove the remaining $O_h(X/y)$ exceptions. Each retained weight is at most one, so positive weighted mass
bounds the cardinality below. This proves \eqref{eq:full-omega-consumer}.
\end{proof}

\bibliographystyle{plainnat}
\bibliography{references}

\section*{Use of AI}
The proofs and the text of this paper were produced with AI systems (GPT-6 Astra, GPT-6.1 Sol and Claude
Opus~5.5) under the author's direction. They were checked by independent reviews with AI models, and the
algebraic identities by computation on finite ranges (programs in the ancillary files). The author takes
responsibility for the content.
\end{document}
